% ============================================================================
% APPENDABLE MANUSCRIPT FILE
% Where the Credit Is: Why Running Counts Teach What Final Counts Miss
% Proposed appendix extension to Theorem 4, prepared 2026-10-13.
%
% USAGE: % ============================================================================
% APPENDABLE MANUSCRIPT FILE
% Where the Credit Is: Why Running Counts Teach What Final Counts Miss
% Proposed appendix extension to Theorem 4, prepared 2026-10-13.
%
% USAGE: % ============================================================================
% APPENDABLE MANUSCRIPT FILE
% Where the Credit Is: Why Running Counts Teach What Final Counts Miss
% Proposed appendix extension to Theorem 4, prepared 2026-10-13.
%
% USAGE: % ============================================================================
% APPENDABLE MANUSCRIPT FILE
% Where the Credit Is: Why Running Counts Teach What Final Counts Miss
% Proposed appendix extension to Theorem 4, prepared 2026-10-13.
%
% USAGE: \input{learned_routing_theorem4_extension.tex} in the manuscript
% appendix. Requires amsmath, amssymb, amsthm (already common in NeurIPS papers).
%
% STATUS: Audited reduced-model extension; quantitative hypotheses are explicit.
% It is an appendix addition and does not alter either part of Theorem 4.
% The end-to-end conclusion is explicitly CONDITIONAL on a symmetry-preserving
% planted-feature parameterization and quantitative parameter-group mobilities.
% It is NOT a theorem for unrestricted random Transformer initialization.
%
% Suggested main-text replacement/addition for Section 4.2 (copy separately):
% "Theorem 4 isolates transition learning with supplied useful directions and,
% in the clean binary extension, learns distractor suppression. A complementary
% full-prefix analysis (Appendix [reference]) considers a weakly seeded
% count--letter feature with no prescribed current-letter position. In a
% symmetry-reduced two-ReLU model, the innovation covariance selects the
% correct relative offset, and fixed-mobility joint gradient flow can learn
% this bias and reach low step loss in polynomial time. This extension still
% relies on tied feature directions, a small parameter-group learning rate,
% and an accelerated trainable relative bias; it does not prove convergence
% for unrestricted Block 2 training or by itself extend Theorem 3's matched
% outcome plateau to the planted activation-pattern initialization."
% ============================================================================

\providecommand{\LRstep}{L_{\mathrm{step}}}
\providecommand{\LRout}{L_{\mathrm{out}}}
\providecommand{\LRsig}{\sigma}
\providecommand{\LRnorm}[1]{\left\lVert #1\right\rVert}
\providecommand{\LRinner}[2]{\left\langle #1,#2\right\rangle}
\theoremstyle{plain}
\newtheorem{lrlemma}[lemma]{Lemma}
\newtheorem{lrproposition}[proposition]{Proposition}
\newtheorem{lrtheorem}[theorem]{Theorem}
\newtheorem{lrremark}[remark]{Remark}
\crefname{lrlemma}{Lemma}{Lemmas}
\crefname{lrproposition}{Proposition}{Propositions}
\crefname{lrtheorem}{Theorem}{Theorems}
\crefname{lrremark}{Remark}{Remarks}

\section{Weakly seeded full-prefix routing with joint feature learning}
\label{app:weak-seeded-full-prefix-routing}

\paragraph{Relation to Theorem 4.}
Theorem~4 of the main paper provides a polynomial-time positive-margin
convergence theorem with prescribed useful attention directions, and a clean
binary variant that learns soft penalties on designated zero-key/value
distractors.  This section studies a different intermediate regime: the model
is supplied a weak \emph{count--letter feature}, but no current-letter position
or preferred relative offset.  Attention remains causal and full-prefix;
relative-position biases are all initially zero.  The purpose is to isolate
which parts of routing and transition learning follow from population
cross-entropy and which require additional architectural restrictions.
Throughout, the terms ``full-prefix'' and ``joint training'' do not mean
unrestricted Transformer parameters: a tied, symmetry-reduced value/MLP
family is maintained as specified below.

\subsection{Data, representation, and the exact initial identities}
\label{app:weak-routing-initial}

Consider binary counting $K=2$, with independent
$b_j\sim\operatorname{Bernoulli}(p)$, $0<p<1/2$, and uniform $s_0$.
The application of interest is $p=1/4$.  As in the main paper, the input
contains the prompt letters and the displayed running counts, and the
prediction at $i_t$ for $s_t$ sees $s_{t-1}$; the associated prompt-letter
site is $j_t$ and
\begin{equation}\label{eq:lr-constant-lag}
 i_t-j_t=n,\qquad t=2,\ldots,n.
\end{equation}
The attention distribution is normalized over the \emph{whole visible causal
prefix}, including query, separator, and displayed-state sites.  We write
$\omega_t=1/(n-1)$ and $\chi=p(1-p)$.

For the actual one-head Block 2 with hidden parameters $\phi$ and readout
$V$, let $h_t=h_\phi(x_t)$, $z_t=Vh_t$, and
\begin{equation}
 q_s=(1-p)e_s+pe_{s\oplus1},\qquad
 g_t=e_{s_t}-q_{s_{t-1}}=(b_t-p)(e_{s_{t-1}\oplus1}-e_{s_{t-1}}).
\end{equation}
The state $s_{t-1}$ is independent of the entire prompt-bit vector: this
follows from uniformity of $s_0$.  Consequently,
\begin{equation}\label{eq:lr-innovation-cov}
 \mathbb E[(b_t-p)(b_j-p)]=\chi\mathbf 1\{j=t\}.
\end{equation}
For fixed $p=1/4$, the right side is a positive constant, unlike the
outcome-correlation factor $\lambda^n$.

\begin{lrlemma}[Exact target decomposition and readout rule credit]
\label{lem:lr-decomposition}
For population step cross-entropy and arbitrary logits,
\begin{equation}\label{eq:lr-kl-decomp}
 \LRstep=H(p)+\mathbb E\operatorname{KL}
 (q_{s_{t-1}}\Vert\operatorname{softmax}(z_t))
 -\Phi_{\mathrm{rule}},\qquad
 \Phi_{\mathrm{rule}}=\mathbb E\LRinner{g_t}{z_t}.
\end{equation}
Writing $C(\phi)=\mathbb E[g_t h_t^\top]$, one has
$\nabla_V\Phi_{\mathrm{rule}}=C(\phi)$.  At an exact state-only predictor
$\operatorname{softmax}(z_t)=q_{s_{t-1}}$, one has
$\nabla_V\LRstep=-C(\phi)$.
\end{lrlemma}
\begin{proof}
The conditional expected label is $q_{s_{t-1}}$ and the centered
innovation is $g_t$.  Expanding cross-entropy as log-sum-exp minus the
target logit gives \eqref{eq:lr-kl-decomp}.  The readout derivatives follow
from $z_t=Vh_t$ and the usual softmax cross-entropy gradient.
\end{proof}

\begin{lrlemma}[Exact zero-readout bootstrap]
\label{lem:lr-zero-bootstrap}
Let $u=(1/2,1/2)$ and
$M(\phi)=\mathbb E[(e_{s_t}-u)h_\phi^\top]$.
Under \emph{unit-mobility Euclidean gradient flow} of $(\phi,V)$ with
$V(0)=0$, at any initial parameter point with a locally continuously differentiable hidden map,
\begin{equation}\label{eq:lr-zero-boot}
 \dot V(0)=M(\phi_0),\qquad \dot\phi(0)=0,\qquad
 \ddot\phi(0)=\nabla_\phi\frac12\LRnorm{M(\phi_0)}_F^2.
\end{equation}
For a group of hidden parameters trained with constant mobility $\nu_G$,
with unit readout mobility, its initial acceleration is $\nu_G$ times the corresponding component of
the final gradient.
\end{lrlemma}
\begin{proof}
At $V=0$, predicted probabilities are $u$, so the readout velocity is
$M$.  The gradient into $\phi$ has a factor $V$ and vanishes initially.
Differentiating the hidden gradient-flow equation at time zero leaves
$(D_\phi M)^\top M$, which is the gradient of $\LRnorm{M}^2/2$.
\end{proof}

For reference, if $P(\phi)=\mathbb E[(q_{s_{t-1}}-u)h_\phi^\top]$ is nonzero,
then the prior-orthogonalized rule feature
\begin{equation}
 R(\phi)=C(\phi)-\frac{\LRinner{C(\phi)}{P(\phi)}_F}
 {\LRnorm{P(\phi)}_F^2}P(\phi)
\end{equation}
satisfies
$\left.\frac{d}{dt}\LRinner{V(t)}{R(\phi(t))}_F\right|_{t=0}
=\LRnorm{R(\phi_0)}_F^2$.
This is an \emph{initial readout-alignment identity}, not a long-time
non-collapse theorem.

\begin{lrlemma}[Exact positional QK geometry]
\label{lem:lr-qk-geometry}
For the reduced positional projections $Q,K$ and
$S=QK^\top/\sqrt{d_h}$, let $G=\nabla_S\LRstep$.
When the positional embeddings are mutually orthogonal vectors
$\varepsilon e_i$ and the underlying positional QK matrix blocks share
mobility $\nu_Q$,
\begin{align}
 \dot Q&=-\frac{\nu_Q\varepsilon^2}{\sqrt{d_h}}GK,
 &\dot K&=-\frac{\nu_Q\varepsilon^2}{\sqrt{d_h}}G^\top Q,\label{eq:lr-qk-flow}\\
 \dot S&=-\frac{\nu_Q\varepsilon^2}{d_h}
 (QQ^\top G+GKK^\top),\label{eq:lr-qk-scoreflow}\\
 \frac{d}{dt}(Q^\top Q-K^\top K)&=0.\label{eq:lr-balanced-conservation}
\end{align}
These identities hold as readout and MLP parameters evolve, provided the
query/key dependence on the positional projections remains as stated.
\end{lrlemma}
\begin{proof}
Apply the chain rule to $QK^\top/\sqrt{d_h}$.  The one-hot positional
features supply the identical mobility factor $\nu_Q\varepsilon^2$ for
each projected row.  Substituting the two velocities in
$\frac{d}{dt}(Q^\top Q-K^\top K)$ cancels the four terms pairwise.
\end{proof}

In particular, if the conserved difference
$Q_0^\top Q_0-K_0^\top K_0$ has at least $n-1$ eigenvalues at least
$\mu>0$, the supervised-query Gram matrix $Q(t)Q(t)^\top$ has its
smallest eigenvalue at least $\mu$ as long as its number of rows is $n-1$.
This is a \emph{conditional} conditioning mechanism, not a conclusion for
arbitrarily initialized full Transformer QKV.

\subsection{A weakly planted feature without a planted route}
\label{app:weak-routing-family}

We now specify the symmetry-reduced family used by the remaining results.
Fix the query letter. Let $s=(-1)^{s_{t-1}}$ and let attention depend on positions, not the
random letter identities.  The full-prefix pooled, centered value is
\begin{equation}\label{eq:lr-centered-value}
 Z_t=\sum_{j\in\mathcal J}a_{tj}(b_j-p),\qquad
 a_{tj}=\frac{\exp(S_{i_tj})}{\sum_{k\le i_t}\exp(S_{i_tk})},
 \quad \mathcal J=\{j_1,\ldots,j_n\}.
\end{equation}
Non-letter positions have zero \emph{value}, not zero attention mass.
A fixed token/position-aware linear value direction realizes the centered contributions
$b_j-p$ for prompt letters and zero contributions on the remaining visible
sites.  The trainable value parameter is its amplitude $v$. For the fixed query, assign its token the usual match
value $1-p$ and subtract $1-p$ using its position embedding; assign count
and separator tokens value zero. Prompt-letter positions have zero positional
correction. Orthogonal token and position registers realize this exactly.
Allowing all
unrestricted value directions to vary is outside this invariant family.

A pair of ReLUs has preactivations
\begin{equation}\label{eq:lr-gates}
 r_+=\operatorname{ReLU}(as+bvZ_t+\beta),\qquad
 r_-=\operatorname{ReLU}(-as+bvZ_t+\beta).
\end{equation}
The binary logit (positive for state sign $+1$) is
\begin{equation}\label{eq:lr-factor-logit}
 f_t=cs+u\omega(r_+-r_-).
\end{equation}
Here $\beta$ is the \emph{MLP bias}, distinct from the attention
relative biases $\beta_\ell^{\mathrm{rel}}$ introduced below.
Every displayed previous state is available directly in the residual stream.

\begin{lrlemma}[Exact two-ReLU feature family]
\label{lem:lr-two-relu}
If
\begin{equation}\label{eq:lr-gate-condition}
 a>|\beta|+|bv|,
\end{equation}
then for all inputs with $|Z_t|\le1$,
\begin{equation}\label{eq:lr-feature-param}
 \begin{aligned}
 r_+-r_-&=s(d+wZ_t),& d&=a+\beta,& w&=bv,\\
 f_t&=s(C+gZ_t),& C&=c+u\omega d,& g&=u\omega w.
 \end{aligned}
\end{equation}
The feature family does not specify which prompt-letter position is current.
\end{lrlemma}
\begin{proof}
When $s=1$, only $r_+$ is active; when $s=-1$, only $r_-$ is active.
The inactive preactivation is strictly negative by
\eqref{eq:lr-gate-condition}.  Substitution gives the stated identity.
\end{proof}

\paragraph{Rule covariance and bounded perturbations.}
For a purely centered hidden coordinate $H_t=\rho sZ_t$,
Equations~\eqref{eq:lr-innovation-cov} and
\eqref{eq:lr-centered-value} give the exact identity
\begin{equation}\label{eq:lr-centered-seed-cov}
 \mathbb E[g_tH_t]=\rho\chi\bar A(e_1-e_0),\qquad
 \mathbb E[(q_{s_{t-1}}-u)H_t]=0.
\end{equation}
In particular the feature covariance is nonzero even when every
letter site has the same initial attention weight.  If a planted seed
has rule covariance at least $\kappa>0$ and the covariance map is
$L$-Lipschitz on a neighbourhood, then every perturbation of norm at
most $\kappa/(2L)$ preserves covariance at least $\kappa/2$.
This elementary robustness statement applies to perturbations within
a parameterization preserving the feature roles; it does not ensure
that arbitrary off-family trainable directions remain stable under
subsequent unrestricted gradient flow.

\paragraph{Initial positional routing acceleration (ideal centered channel).}
If the hidden representation contains the \emph{exact} independent
coordinate $H_t=\rho sZ_t$, all its other target-correlated coordinates
are independent of positional QK, and $V(0)=0$, then
$F(\phi)=\LRnorm{M(\phi)}_F^2/2$ satisfies, upon varying positional QK,
$F=F_0+\rho^2\chi^2\bar A^2$.  Thus under unit QK mobility,
\begin{equation}\label{eq:lr-centered-accel}
 \dot{\bar A}(0)=0,\qquad
 \ddot{\bar A}(0)=2\rho^2\chi^2\bar A(0)
 \LRnorm{\nabla_{Q,K}\bar A(0)}_F^2\ge0.
\end{equation}
Nondegenerate positional QK gives strict initial acceleration.
This statement is \emph{conditional on exact feature orthogonalization}:
the raw two-ReLU coordinate
$s(d+wZ_t)$ also contains prior information, which creates additional
terms.  It is therefore not a general acceleration theorem for arbitrary
planted Block 2 parameters.

The two hidden units may be trained as an \emph{explicitly tied Euclidean
submanifold}: opposite count weights, equal centered-letter weights and
biases, and opposite output weights.  The induced Euclidean metric assigns
factor $1/2$ to a gradient update of a scalar appearing twice with equal
magnitude.  This tied family is a model restriction; symmetry alone does
not prove its invariance for all unrestricted Transformer parameters.
The mobilities below define flow in the displayed effective factor coordinates.
If a matrix realization uses a fixed direction of nonunit norm, its norm
factor must be absorbed into the assigned coordinate mobility; the scalar
equations are not a claim about unnormalized full-matrix Euclidean flow.

Take $u(0)=c(0)=0$, $\omega_0>0$, $b_0<0$, $v_0>0$, and
$a_0>|\beta_0|+|b_0v_0|$.  Assume all planted magnitudes and gate margins
are inverse-polynomially bounded.  Initialize the positional attention
scores without a preferred correct position and set all trainable
relative-position biases to zero.  This is a \emph{weak computational seed},
not random discovery of the count--letter interaction.  Unlike the
fixed-sign-ReLU-bias initialization of Assumption~2 in the main paper,
the gates here deliberately respond differently to the displayed count.
Thus Theorem~3 does not automatically give a matched outcome plateau for
this precise planted initialization.

\subsection{Exact cross-entropy routing signs}
\label{app:weak-routing-signs}

Let $\sigma(x)=(1+e^{-x})^{-1}$, $y_t'=1-b_t$ and
\begin{equation}\label{eq:lr-logistic-risk}
 \LRstep=\sum_{t=2}^n\omega_t\,\mathbb E
 \left[\log(1+e^{C+gZ_t})-y_t'(C+gZ_t)\right].
\end{equation}
This is exactly the binary cross-entropy of
\eqref{eq:lr-feature-param} because multiplying the binary logit and
label by $s$ cancels the previous-state sign.
Let $A_t=a_{t,j_t}$ and $\bar A=\sum_t\omega_t A_t$.

\begin{lrlemma}[Innovation selects the current source]
\label{lem:lr-signs}
For position-only attention,
\begin{equation}\label{eq:lr-initial-gate-gradient}
 \left.\frac{\partial\LRstep}{\partial g}\right|_{g=0}
 =\chi\bar A>0
\end{equation}
for every $C$.  At any finite logits and $g<0$, define
$v_{tj}=\partial L_t/\partial a_{tj}$, where $L_t$ is the single-time
population risk without $\omega_t$.  Then
\begin{equation}\label{eq:lr-feature-source-signs}
 v_{t,j_t}<0,\qquad v_{tj}\ge0
 \quad(j\in\mathcal J,\;j\ne j_t),\qquad
 v_{tj}=0\quad(j\notin\mathcal J).
\end{equation}
If $|C|+|g|\le B$, then
\begin{equation}\label{eq:lr-source-bounds}
 v_{t,j_t}\le-2\chi|g|\sigma(-B),\qquad
 0\le v_{tj}\le\frac{\chi g^2a_{tj}}4\quad(j\ne j_t\text{ a letter}).
\end{equation}
\end{lrlemma}
\begin{proof}
Centering and independence give
$\mathbb E[(b_t-p)Z_t]=\chi A_t$, yielding
\eqref{eq:lr-initial-gate-gradient}.  For the correct letter,
conditioning on the other bits gives
\[
 v_{t,j_t}=g\chi\mathbb E
 \{1+\sigma(U+gA_t(1-p))-\sigma(U-gA_tp)\},
\]
where $U$ is the contribution of all other positions and $C$.
The braced term is strictly positive.  Under the logit bound it is at
least $2\sigma(-B)$.  For an incorrect letter $j$, the target is
independent of $b_j$, and conditioning on all other bits yields
\[
 v_{tj}=g\chi\mathbb E
 \{\sigma(U+ga_{tj}(1-p))-\sigma(U-ga_{tj}p)\}.
\]
For $g<0$ the two terms in braces are nonpositive; their magnitude is
at most $|g|a_{tj}/4$ by $\sigma'\le1/4$.  Non-letter values are zero.
\end{proof}

Put $G_{tj}=\omega_ta_{tj}(v_{tj}-\bar v_t)$, with
$\bar v_t=\sum_ka_{tk}v_{tk}$; this is the gradient with respect to a
pre-softmax score.  The following global comparison will be important.

\begin{lrlemma}[Correct-versus-incorrect score comparison]
\label{lem:lr-postrouting-gradient}
Whenever $g<0$ and $A_t\ge1/2$,
\begin{equation}\label{eq:lr-score-comparison}
 G_{t,j_t}<0,\qquad G_{tj}-G_{t,j_t}\ge0
 \quad\text{for every other visible site }j.
\end{equation}
Consequently, if the attention biases are shared by relative offset and
$i_t-j_t=n$, then for every $\ell\ne n$,
\begin{equation}\label{eq:lr-bias-gap-mono}
 \frac{d}{dt}
 (\beta_n^{\mathrm{rel}}-\beta_\ell^{\mathrm{rel}})\ge0
\end{equation}
under common positive relative-bias mobility.
\end{lrlemma}
\begin{proof}
The first inequality follows because the correct attention-weight
derivative is strictly smaller than all others, so it lies below their
weighted mean.  For an incorrect site with $a=a_{tj}\le A=A_t$ and
$v=v_{tj}\ge0$, $v_c=v_{t,j_t}<0$,
\[
 \frac{G_{tj}-G_{t,j_t}}{\omega_t}
 =av-Av_c+(A-a)\bar v_t
 \ge av+A(1-A+a)(-v_c)\ge0.
\]
We used $\bar v_t\ge Av_c$ and $A-a\ge0$.
An offset $\ell$ occurs at most once in any attention row.  The correct
offset $n$ appears in every supervised row.  Summing the comparisons
within rows where $\ell$ occurs, and using $G_{t,j_t}<0$ in the
remaining rows, gives
$\partial_{\beta_n^{\mathrm{rel}}}\LRstep
\le\partial_{\beta_\ell^{\mathrm{rel}}}\LRstep$.
\end{proof}

\subsection{Feature bootstrap with every parameter group active}
\label{app:weak-routing-bootstrap}

All training uses one \emph{fixed} system of positive mobilities:
$c,u$ have mobility one; tied MLP/value factors have mobility
$\nu_F>0$; positional QK have mobility $\nu_Q>0$; all relative biases
have common mobility $\nu_R>0$.
No parameter group is explicitly frozen or switched at a chosen time.
The factorized feature family and two-ReLU gate condition are nevertheless
restrictions on the parameterization.

Write $\alpha=u\omega$, $w=bv<0$, $d=a+\beta>0$,
$L_C=\partial\LRstep/\partial C$, $L_g=\partial\LRstep/\partial g$.
While the gate regime holds, the tied Euclidean gradient equations include
\begin{align}
 \dot c&=-L_C,
 &\dot u&=-\omega(dL_C+wL_g),\label{eq:lr-factor-u}\\
 \dot\omega&=-\tfrac{\nu_Fu}{2}(dL_C+wL_g),
 &\dot a=\dot\beta&=-\tfrac{\nu_F\alpha}{2}L_C,\\
 \dot b&=-\tfrac{\nu_F\alpha v}{2}L_g,
 &\dot v&=-\nu_F\alpha bL_g.
 \label{eq:lr-factor-bv}
\end{align}
In particular,
\begin{equation}\label{eq:lr-factor-invariants}
 u^2-\frac{2}{\nu_F}\omega^2=-\frac{2}{\nu_F}\omega_0^2,
 \qquad
 v^2-2b^2=v_0^2-2b_0^2,
 \qquad a-\beta=a_0-\beta_0.
\end{equation}
The first invariant gives $\omega^2\ge\omega_0^2$.  Under the gate
condition, the independent prior readout satisfies
$c(t)\le c_\star=\log((1-p)/p)$: at $c=c_\star$, with $u\omega\ge0$,
its derivative points inward.  At $u=0$ the velocity points into $u>0$,
and $w=bv$ retains its negative sign as long as the bounded gate regime
holds: on either boundary $b=0$ or $v=0$, one has $L_g=\chi\bar A>0$,
so the vector field points into $b<0,v>0$ (the simultaneous zero is excluded
by uniqueness from the nonzero seed). Thus $g=u\omega w<0$ for positive times.

\begin{lrlemma}[Quantitative weak-feature bootstrap, with explicit exits]
\label{lem:lr-feature-bootstrap}
Suppose on an interval prior to reaching the target magnitude that
\begin{equation}\label{eq:lr-amplification-region}
 \bar A\ge A_{\min}>0,\quad
 d\le D_F,\quad
 0<W\le|w|\le W_{\max},\quad
 \omega\ge\omega_{\min}>0,
\end{equation}
and the gate condition \eqref{eq:lr-gate-condition} holds.  Put
\begin{equation}\label{eq:lr-gamma-bootstrap}
 H=D_F+W_{\max},\quad
 \gamma=\chi A_{\min}\left(\frac W H\right)^2,
 \quad
 T_{\mathrm{boot}}=\frac4{3\omega_{\min}^2H^2}.
\end{equation}
Starting with $c=u=0$, the quantity $x=-g$ reaches $x\ge\gamma$
by $T_{\mathrm{boot}}$, unless one of the region conditions fails first.
No assumption $L_C\le0$ is required.
\end{lrlemma}
\begin{proof}
At $c\le c_\star$, let $\alpha=u\omega$.  Since
$0<d+wZ_t\le H$, logistic curvature implies
$L_C\le\alpha H/4$.  At $g=0$, $L_g=\chi\bar A$ and the
same curvature bound gives $L_g\ge\chi A_{\min}-x/4$.
When $x\le\gamma$,
\[
 dL_C+wL_g
 \le\frac{D_FH}{4W}x-W\chi A_{\min}
 +\frac{W_{\max}}4x
 \le-\frac34W\chi A_{\min}.
\]
The last step uses
$D_FH+WW_{\max}\le H^2$.
From \eqref{eq:lr-factor-u} and
$\omega^2\ge\omega_{\min}^2$,
$\dot\alpha\ge(3/4)\omega_{\min}^2W\chi A_{\min}$.
In this region $L_g>0$ and
$\dot w=-\nu_F\alpha(v^2+b^2/2)L_g\le0$.
Hence $\dot x\ge(3/4)\omega_{\min}^2W^2\chi A_{\min}$,
which gives the claimed hitting time.
\end{proof}

The groupwise loss budget is exact in the induced Euclidean metrics:
for the tied feature coordinates $(a,\beta,b,\omega,v)$, the squared
metric weights are $(2,2,2,2,1)$. For projected positional rows $(Q,K)$,
the squared metric weight is $\varepsilon^{-2}$.
\begin{equation}\label{eq:lr-mobility-budget}
 -\dot{\LRstep}=
 \LRnorm{\nabla_{c,u}\LRstep}^2
 +\nu_F\LRnorm{\nabla_F\LRstep}^2
 +\nu_Q\LRnorm{\nabla_{Q,K}\LRstep}^2
 +\nu_R\LRnorm{\nabla_{\beta^{\mathrm{rel}}}\LRstep}^2.
\end{equation}
Here group norms and gradients use those induced metrics. For each group
with mobility $\nu$, Cauchy--Schwarz implies
\begin{equation}\label{eq:lr-movement-budget}
 \LRnorm{\theta_G(T)-\theta_G(0)}
 \le\sqrt{\nu T\,[\LRstep(0)-\LRstep(T)]}
 \le\sqrt{\nu T\log2}.
\end{equation}
Thus sufficiently small but positive $\nu_F,\nu_Q$ preserve prescribed
strict feature-gate and base-score margins over any known finite horizon.
This justifies a controlled bootstrap; it does not establish unit-rate
stability in the unrestricted architecture.

\subsection{Relative-position learning and post-routing convergence}
\label{app:weak-routing-convergence}

Use positional scores
\begin{equation}\label{eq:lr-relative-scores}
 S_{i_tj}=B_{i_tj}+\beta_{i_t-j}^{\mathrm{rel}},
 \qquad
 B_{i_tj}=\frac{q_{i_t}^\top k_j}{\sqrt{d_h}},
 \qquad \beta_\ell^{\mathrm{rel}}(0)=0.
\end{equation}
The correct offset $n$ is not privileged at initialization.

\begin{lrlemma}[Correct-offset learning in the small-feature region]
\label{lem:lr-bias-learning}
Assume all base scores satisfy $|B_{i_tj}|\le D_A$,
all relative biases start at zero, and
$0<x=-g\le A_{\min}/4$. If $A_t\ge A_{\min}>0$,
then the correct score derivative is negative and all other visible-site
score derivatives are nonnegative:
\begin{equation}\label{eq:lr-smallscore-sign}
 G_{t,j_t}<0,\qquad G_{tj}\ge0\ (j\ne j_t).
\end{equation}
Consequently, for the shared relative biases,
\begin{equation}\label{eq:lr-beta-routes}
 \dot\beta_n^{\mathrm{rel}}>0,\qquad
 \dot\beta_\ell^{\mathrm{rel}}\le0\quad(\ell\ne n).
\end{equation}
They imply the uniform floor
\begin{equation}\label{eq:lr-floor}
 A_t\ge A_{\min}:=\frac{e^{-2D_A}}m,
 \qquad m:=\max_t\#\{j:j\le i_t\}\le3n.
\end{equation}
If the initial bootstrap reaches $x=\gamma$ and
$\gamma\le A_{\min}/4$, there is a polynomially bounded fixed bias
mobility $\nu_R$ that drives every $A_t$ to $2/3$ before $x$ leaves
$[\gamma/2,A_{\min}/4]$, provided the feature and base-score
neighbourhoods hold over that interval.
\end{lrlemma}
\begin{proof}
By Lemma~\ref{lem:lr-signs}, the correct value derivative is negative;
incorrect-letter derivatives are nonnegative.  The correct derivative has the logit-independent bound
\[
 -v_{t,j_t}\ge\chi x(1-xA_t/4)\ge\chi x(1-x/4),
\]
because $\sigma$ is $1/4$-Lipschitz. Hence
\[
 \bar v_t\le-\chi x A_t(1-x/4)+\chi x^2/4<0
\]
when $x\le A_{\min}/4$ and $A_t\ge A_{\min}$.  Hence the softmax derivatives have the stated signs.
The bias sign conclusion follows from the fixed offset
\eqref{eq:lr-constant-lag}.
With $\beta_n^{\mathrm{rel}}\ge0$, all other biases $\le0$,
and $|B|\le D_A$, softmax normalization yields
$A_t\ge e^{-2D_A}/m$.

For the time estimate, whenever some row has $A_t\le2/3$, the exact
correct-score derivative satisfies
$-G_{t,j_t}\ge c_0\omega_t\chi xA_{\min}$
with $c_0=1/4$ in this small-$x$ region: indeed
$-G_{t,j_t}\ge\omega_t A_t(1-A_t)(-v_{t,j_t})$.
Therefore
$\dot\beta_n^{\mathrm{rel}}
\ge c_0\nu_R\chi\gamma A_{\min}/[2(n-1)]$
as long as $x\ge\gamma/2$ and routing is incomplete.
All rows achieve $A_t\ge2/3$ when
$\beta_n^{\mathrm{rel}}\ge
2D_A+\log(2(m-1))$.
Thus the additional time is at most
\[
 T_{\mathrm{route}}
 \le \frac{2(n-1)[2D_A+\log(2(m-1))]}
 {c_0\nu_R\chi\gamma A_{\min}}.
\]
On a bounded feature neighbourhood, $|\dot x|\le M_x$ with $M_x$
independent of $\nu_R$ (attention probabilities stay in $[0,1]$).
Choose the fixed $\nu_R$ large enough that this time is less than
$\min\{1,\gamma/(4M_x),(A_{\min}/4-\gamma)/(2M_x)\}$.
For $p=1/4$, the choice $\gamma=\chi A_{\min}(W/H)^2$
automatically gives $\gamma\le3A_{\min}/16<A_{\min}/4$.
No learning-rate switch is required.
\end{proof}

\begin{lrlemma}[Post-routing invariance under bounded QK-score drift]
\label{lem:lr-post-routing}
Suppose $g<0$, $A_t(t_0)\ge2/3$ for all supervised $t$, and the
base-score contrasts change by at most $\log(5/4)$ after $t_0$:
\[
 |[B_{i_t,j_t}-B_{i_tj}](t)-[B_{i_t,j_t}-B_{i_tj}](t_0)|
 \le\log(5/4).
\]
As long as $g<0$ and the centered-feature family holds, all later
correct-position attention weights obey
\begin{equation}\label{eq:lr-post-floor}
 A_t(t)\ge\frac8{13}>\frac35.
\end{equation}
\end{lrlemma}
\begin{proof}
As long as all $A_t\ge1/2$, the bias comparisons of
Lemma~\ref{lem:lr-postrouting-gradient} imply that each relative-bias
gap $\beta_n^{\mathrm{rel}}-\beta_\ell^{\mathrm{rel}}$ is
nondecreasing.  The assumed bound on the other score contrasts says
the odds of each competitor relative to the correct site can increase
by at most a factor $5/4$.  Since the sum of competitor odds is at most
$1/2$ at time $t_0$, it stays at most $5/8$, yielding
$A_t\ge1/(1+5/8)=8/13$.  A first-exit argument closes the condition
$A_t\ge1/2$.
\end{proof}

\begin{lrlemma}[Uniform separability and convex-readout descent]
\label{lem:lr-readout-convergence}
Suppose $A_t\ge3/5$ at every supervised position.  Then
\begin{equation}\label{eq:lr-bit-separation}
 (1-2b_t)\left(\frac12-p-Z_t\right)\ge\frac1{10}
\end{equation}
for every prompt-bit realization.  Suppose in addition that, on a
finite interval after $t_0$, the trainable feature factors remain in a
neighbourhood that admits a \emph{fixed} readout comparator
$\widehat r=(\widehat c,\widehat u)$ satisfying
\begin{equation}\label{eq:lr-comparator}
 \LRstep(\widehat r;\phi(t))\le\eta/2
 \quad\text{for all }t\text{ in the interval}.
\end{equation}
If the two readout coefficients $(c,u)$ continue to train at unit
mobility and every other group has nonnegative mobility, then
\begin{equation}\label{eq:lr-descent}
 \LRstep(t_0+T)\le\frac\eta2+
 \frac{\LRnorm{r(t_0)-\widehat r}^2}{2T}.
\end{equation}
In particular, $T=\LRnorm{r(t_0)-\widehat r}^2/\eta$ suffices for
$\LRstep\le\eta$.
\end{lrlemma}
\begin{proof}
For $b_t=0$, $Z_t\le1-p-A_t$; for $b_t=1$,
$Z_t\ge A_t-p$.  Both cases imply
\eqref{eq:lr-bit-separation}.
For a reference seed $(d_0,w_0,\omega_0)$ with $w_0<0$, take
\[
 R=20\log(2/\eta),\qquad
 \widehat u=\frac{R}{\omega_0|w_0|},\qquad
 \widehat c=R\left(\frac12-p-\frac{d_0}{|w_0|}\right).
\]
At the reference feature, its effective logit is
$R(1/2-p-Z_t)$, which has signed margin at least $R/10$.
Require the feature neighbourhood to satisfy
\begin{equation}\label{eq:lr-comparator-stability}
 |\omega d-\omega_0d_0|+|\omega w-\omega_0w_0|
 \le\frac{\omega_0|w_0|}{20}.
\end{equation}
Since $|Z_t|\le1$, the logit perturbation is at most $R/20$.
The signed margin remains at least $R/20$, so the comparator loss is
at most $e^{-R/20}=\eta/2$. The same comparator is fixed before training,
independently of the routing hitting time.
Cross-entropy is convex in $r=(c,u)$ at fixed feature factors.  Thus
$\frac12\frac{d}{dt}\LRnorm{r-\widehat r}^2
\le\LRstep(\widehat r;\phi(t))-\LRstep(r;\phi(t))$.
The full joint flow makes the loss nonincreasing.  Integrating gives
\eqref{eq:lr-descent}.
\end{proof}

\begin{lrtheorem}[Finite-accuracy learned routing in a symmetry-reduced family]
\label{thm:lr-conditional-bridge}
Fix $n\ge2$, $p=1/4$, $0<\eta<1/2$ and $0<\varepsilon\le1$.
Use the explicitly tied family \eqref{eq:lr-gates}--\eqref{eq:lr-feature-param}
with $c_0=u_0=0$, $b_0<0$, $v_0,\omega_0>0$, and
$a_0>|\beta_0|+|b_0v_0|$. The fixed query and centered value direction
are as in \cref{app:weak-routing-family}. QK consists only of trainable
positional projections of orthogonal position embeddings; every visible
causal site participates in softmax. All relative biases start at zero.

Assume a complexity bound $P\ge2$ bounds $n$, the widths,
$\varepsilon^{-1}$, all initial factor magnitudes, and the Frobenius norms
of the projected QK matrices. Also require
\[
 \min\{|b_0|,v_0,\omega_0,a_0-|\beta_0|-|b_0v_0|\}\ge P^{-1}.
\]
The initial base scores must have moderate spread: for some $D_0\ge0$,
\[
 |B_{i_tj}(0)|\le D_0,\qquad e^{2D_0}\le P.
\]
No correct relative offset is privileged by the relative biases. The base
scores may be bounded perturbations of equal scores; the assertion does
not require an initially favorable correct-site score.

There exist fixed positive mobilities $\nu_F,\nu_Q$ and $\nu_R$, with
$\nu_F^{-1},\nu_Q^{-1},\nu_R$ bounded by polynomials in $P$ and $1/\eta$,
such that the uninterrupted joint population gradient flow reaches
\begin{equation}\label{eq:lr-finite-theorem}
 \LRstep(T)\le\eta
\end{equation}
for $T\le\operatorname{poly}(P,1/\eta)$.
Before the post-routing convergence interval it reaches $A_t\ge2/3$
for every supervised $t$, and preserves $A_t\ge8/13$ throughout that
interval. The mobilities are chosen once at initialization, as below.

The claim concerns a \emph{target-dependent finite-accuracy schedule of
constant group mobilities}, not asymptotic loss convergence under one
fixed mobility choice independent of $\eta$, and not the unrestricted
Transformer parameterization.
\end{lrtheorem}
\begin{proof}[Proof with mobilities fixed before training]
Choose feature and QK balls of radii $r_F,r_Q>0$ around their initial
coordinates. Their inverse radii can be polynomially bounded in $P$.
For the feature ball require the gate condition, $b<0$, $v,\omega>0$,
\eqref{eq:lr-comparator-stability}, and uniform bounds
\[
 0<W\le|w|\le W_{\max},\quad d\le D_F,\quad
 \omega_{\min}\le\omega\le\Omega.
\]
The quantities $W^{-1},\omega_{\min}^{-1},W_{\max},D_F,\Omega$ are
polynomial in $P$. For the QK ball require $|B|\le D_A=D_0+1$ and
\[
 |(B_{i_t,j_t}-B_{i_tj})-(B_{i_t,j_t}(0)-B_{i_tj}(0))|
 \le\tfrac12\log(5/4).
\]
Polynomial maps on bounded seed balls are Lipschitz with polynomial
constants; the assumed inverse-polynomial gate and factor margins give
these radii. Both QK points at any two times in the ball therefore have
contrast differences at most $\log(5/4)$.

Set $A_{\min}=e^{-2D_A}/m$, $H=D_F+W_{\max}$,
$\gamma=\chi A_{\min}(W/H)^2$ and $T_b=T_{\mathrm{boot}}$ from
\eqref{eq:lr-gamma-bootstrap}. These quantities and their relevant
inverses are polynomial in $P$, and $\gamma\le3A_{\min}/16$.
Set $T_{\mathrm{pre}}=T_b+1$. Until that horizon, inside the feature
ball, the unit-mobility readout equations give
\[
 |\dot c|\le1,\qquad |\dot u|\le\Omega H,
 \qquad \|(c,u)\|\le T_{\mathrm{pre}}\sqrt{1+\Omega^2H^2}.
\]
These bounds use only $|\sigma-y'|\le1$ and $|Z_t|\le1$ and do not
depend on bias mobility. With $\nu_F\le1$, the factor equations also
give a polynomial bound $M_x\ge1$ on $|\dot x|$ in this interval,
uniformly over every attention distribution. Choose
\[
 \tau=\min\{1,\gamma/(4M_x),(A_{\min}/4-\gamma)/(2M_x)\},
 \quad K_R=2D_A+\log(2(m-1)),
\]
\[
 \nu_R=\max\left\{1,
 \frac{2(n-1)K_R}{c_0\chi\gamma A_{\min}\tau}\right\},
 \qquad c_0=1/4.
\]
This fixes the relative-bias mobility in advance and makes the routing
bound at most $\tau$. It is polynomially bounded.

Use the fixed reference comparator from the readout lemma and define
\[
 D=T_{\mathrm{pre}}\sqrt{1+\Omega^2H^2}+\|\widehat r\|,
 \qquad T_{\mathrm{post}}=D^2/\eta,\qquad
 T_*=T_{\mathrm{pre}}+T_{\mathrm{post}}.
\]
Now fix the other mobilities by
\[
 \nu_F=\min\{1,r_F^2/(4T_*\log2)\},\qquad
 \nu_Q=\min\{1,r_Q^2/(4T_*\log2)\}.
\]
Their inverses are polynomially bounded. For the tied feature metric,
all squared weights are at least one; for projected QK the weight is
$\varepsilon^{-2}\ge1$. Thus the movement budget bounds ordinary
coordinate movement by $r_F/2$ and $r_Q/2$ up to $T_*$.
A first-exit argument therefore keeps both groups strictly inside their
chosen balls for the whole argument. Finite-time existence follows from
the smooth loss and its gradient-flow energy budget: all mobilities are
finite positive constants, so finite-time parameter escape is excluded.

Inside these balls, $c\le c_\star$ and $u\ge0$ persist by the inward
boundary signs stated above; $\omega>0$ and $w<0$ give $g<0$ for positive
times. In the small-feature region the bias signs imply the attention
floor by a first-exit argument from zero biases. The bootstrap lemma
reaches $x=\gamma$ by $T_b$ unless all routing is already complete;
even in that case the same bootstrap bound applies. Starting at this
hitting time, within at most $\tau$ the bias lemma makes every
$A_t\ge2/3$ (if this already holds, the extra time is zero). Let $t_0$
be the end of this interval; then $t_0\le T_{\mathrm{pre}}$.
The derivative bound keeps $x$ inside $[\gamma/2,A_{\min}/4]$ during
this routing interval. The bounded base-score contrasts and the
post-routing lemma preserve $A_t\ge8/13>3/5$ afterwards.
The feature ball preserves the fixed comparator for the entire remaining
interval. Its initial distance at $t_0$ is at most $D$, so the readout
lemma gives $\LRstep(t_0+T_{\mathrm{post}})\le\eta$.
All mobilities were selected from initial uniform ball bounds; none
requires a later switch or a posterior choice of $t_0$.
\end{proof}

\subsection{Limitations and relationship to the main comparison}
\label{app:weak-routing-limitations}

\paragraph{What is new.}
The statistical signal for the correct relative offset is derived from
independent innovations rather than inserted in the attention mask or bias
initialization.  All visible causal sites compete in the softmax.  The
tied two-ReLU coefficients, value amplitude, positional QK, linear readout,
and relative-position biases train \emph{simultaneously}, although at
different fixed mobilities.  The constant-offset identity
\eqref{eq:lr-constant-lag} is a genuine architectural inductive bias.

\paragraph{What is not proved.}
The reduced family preselects a state--letter feature direction and
centers prompt-letter values.  It excludes unrestricted content-QKV
updates and arbitrary perturbations that leave this tied family.
The result also permits a large relative-bias mobility and very small
positive feature/QK mobilities, and both may depend on the target error
$\eta$.  This is therefore weaker than a theorem for ordinary
unit-mobility, randomly initialized Block 2.  In addition, the
count-dependent ReLU activity differs from the fixed-sign near-zero
bias regime in the main paper's Assumption~2.  The original outcome
plateau (Theorem~3) cannot be transferred to this planted seed without
an additional argument.  Nor does the finite-accuracy guarantee imply
$\LRstep(t)\to0$ for a single mobility choice independent of $\eta$.

\paragraph{Scope of the audited parameterization.}
For interpreting this appendix extension, the tied-subfamily
symmetry, parameter scaling relative to Assumption~2, and value-role restriction must be distinguished
from unrestricted updates in the original computational graph. The quantitative construction above
uses the stated reduced-coordinate metrics, moderate base-score spread,
and bounded seeds; it does not transfer to arbitrary parameter metrics.  In particular, the
centering of the fixed query-letter token must be implemented by the
value projection rather than ignored, the ordinary token-content QK
directions must be either controlled or explicitly excluded, and any
extension to Theorem~3 must separately verify its activation-pattern
hypotheses.  An earlier inference that \emph{nonzero rule covariance by
itself} forces correct routing is false; count-prior interference can
make the initial positional routing acceleration negative.

\paragraph{Suggested interpretation.}
This result is a complementary \emph{conditional routing bridge} next
to the original Theorem~4, not a replacement of its rigorous
positive-margin/global-loss convergence statement.  The next genuinely
stronger theorem would remove the tied feature restriction while
keeping a planted feature, and then establish feature stability,
relative-position learning, and polynomial-time low loss for all
trainable Block 2 directions, with a matched outcome comparison.

\paragraph{Numerical audit.} Exact Bernoulli enumeration and automatic differentiation in
\nolinkurl{handcoded/theory_checks/learned_routing_extension_checks.py}
check the source gradients, small-feature and post-routing score signs, tied-factor
invariants, positional QK metric, uniform separability, and the fixed comparator.
These checks support the displayed local identities and inequalities; they are
not a formal verification of the global theorem or a test of unrestricted training.
 in the manuscript
% appendix. Requires amsmath, amssymb, amsthm (already common in NeurIPS papers).
%
% STATUS: Audited reduced-model extension; quantitative hypotheses are explicit.
% It is an appendix addition and does not alter either part of Theorem 4.
% The end-to-end conclusion is explicitly CONDITIONAL on a symmetry-preserving
% planted-feature parameterization and quantitative parameter-group mobilities.
% It is NOT a theorem for unrestricted random Transformer initialization.
%
% Suggested main-text replacement/addition for Section 4.2 (copy separately):
% "Theorem 4 isolates transition learning with supplied useful directions and,
% in the clean binary extension, learns distractor suppression. A complementary
% full-prefix analysis (Appendix [reference]) considers a weakly seeded
% count--letter feature with no prescribed current-letter position. In a
% symmetry-reduced two-ReLU model, the innovation covariance selects the
% correct relative offset, and fixed-mobility joint gradient flow can learn
% this bias and reach low step loss in polynomial time. This extension still
% relies on tied feature directions, a small parameter-group learning rate,
% and an accelerated trainable relative bias; it does not prove convergence
% for unrestricted Block 2 training or by itself extend Theorem 3's matched
% outcome plateau to the planted activation-pattern initialization."
% ============================================================================

\providecommand{\LRstep}{L_{\mathrm{step}}}
\providecommand{\LRout}{L_{\mathrm{out}}}
\providecommand{\LRsig}{\sigma}
\providecommand{\LRnorm}[1]{\left\lVert #1\right\rVert}
\providecommand{\LRinner}[2]{\left\langle #1,#2\right\rangle}
\theoremstyle{plain}
\newtheorem{lrlemma}[lemma]{Lemma}
\newtheorem{lrproposition}[proposition]{Proposition}
\newtheorem{lrtheorem}[theorem]{Theorem}
\newtheorem{lrremark}[remark]{Remark}
\crefname{lrlemma}{Lemma}{Lemmas}
\crefname{lrproposition}{Proposition}{Propositions}
\crefname{lrtheorem}{Theorem}{Theorems}
\crefname{lrremark}{Remark}{Remarks}

\section{Weakly seeded full-prefix routing with joint feature learning}
\label{app:weak-seeded-full-prefix-routing}

\paragraph{Relation to Theorem 4.}
Theorem~4 of the main paper provides a polynomial-time positive-margin
convergence theorem with prescribed useful attention directions, and a clean
binary variant that learns soft penalties on designated zero-key/value
distractors.  This section studies a different intermediate regime: the model
is supplied a weak \emph{count--letter feature}, but no current-letter position
or preferred relative offset.  Attention remains causal and full-prefix;
relative-position biases are all initially zero.  The purpose is to isolate
which parts of routing and transition learning follow from population
cross-entropy and which require additional architectural restrictions.
Throughout, the terms ``full-prefix'' and ``joint training'' do not mean
unrestricted Transformer parameters: a tied, symmetry-reduced value/MLP
family is maintained as specified below.

\subsection{Data, representation, and the exact initial identities}
\label{app:weak-routing-initial}

Consider binary counting $K=2$, with independent
$b_j\sim\operatorname{Bernoulli}(p)$, $0<p<1/2$, and uniform $s_0$.
The application of interest is $p=1/4$.  As in the main paper, the input
contains the prompt letters and the displayed running counts, and the
prediction at $i_t$ for $s_t$ sees $s_{t-1}$; the associated prompt-letter
site is $j_t$ and
\begin{equation}\label{eq:lr-constant-lag}
 i_t-j_t=n,\qquad t=2,\ldots,n.
\end{equation}
The attention distribution is normalized over the \emph{whole visible causal
prefix}, including query, separator, and displayed-state sites.  We write
$\omega_t=1/(n-1)$ and $\chi=p(1-p)$.

For the actual one-head Block 2 with hidden parameters $\phi$ and readout
$V$, let $h_t=h_\phi(x_t)$, $z_t=Vh_t$, and
\begin{equation}
 q_s=(1-p)e_s+pe_{s\oplus1},\qquad
 g_t=e_{s_t}-q_{s_{t-1}}=(b_t-p)(e_{s_{t-1}\oplus1}-e_{s_{t-1}}).
\end{equation}
The state $s_{t-1}$ is independent of the entire prompt-bit vector: this
follows from uniformity of $s_0$.  Consequently,
\begin{equation}\label{eq:lr-innovation-cov}
 \mathbb E[(b_t-p)(b_j-p)]=\chi\mathbf 1\{j=t\}.
\end{equation}
For fixed $p=1/4$, the right side is a positive constant, unlike the
outcome-correlation factor $\lambda^n$.

\begin{lrlemma}[Exact target decomposition and readout rule credit]
\label{lem:lr-decomposition}
For population step cross-entropy and arbitrary logits,
\begin{equation}\label{eq:lr-kl-decomp}
 \LRstep=H(p)+\mathbb E\operatorname{KL}
 (q_{s_{t-1}}\Vert\operatorname{softmax}(z_t))
 -\Phi_{\mathrm{rule}},\qquad
 \Phi_{\mathrm{rule}}=\mathbb E\LRinner{g_t}{z_t}.
\end{equation}
Writing $C(\phi)=\mathbb E[g_t h_t^\top]$, one has
$\nabla_V\Phi_{\mathrm{rule}}=C(\phi)$.  At an exact state-only predictor
$\operatorname{softmax}(z_t)=q_{s_{t-1}}$, one has
$\nabla_V\LRstep=-C(\phi)$.
\end{lrlemma}
\begin{proof}
The conditional expected label is $q_{s_{t-1}}$ and the centered
innovation is $g_t$.  Expanding cross-entropy as log-sum-exp minus the
target logit gives \eqref{eq:lr-kl-decomp}.  The readout derivatives follow
from $z_t=Vh_t$ and the usual softmax cross-entropy gradient.
\end{proof}

\begin{lrlemma}[Exact zero-readout bootstrap]
\label{lem:lr-zero-bootstrap}
Let $u=(1/2,1/2)$ and
$M(\phi)=\mathbb E[(e_{s_t}-u)h_\phi^\top]$.
Under \emph{unit-mobility Euclidean gradient flow} of $(\phi,V)$ with
$V(0)=0$, at any initial parameter point with a locally continuously differentiable hidden map,
\begin{equation}\label{eq:lr-zero-boot}
 \dot V(0)=M(\phi_0),\qquad \dot\phi(0)=0,\qquad
 \ddot\phi(0)=\nabla_\phi\frac12\LRnorm{M(\phi_0)}_F^2.
\end{equation}
For a group of hidden parameters trained with constant mobility $\nu_G$,
with unit readout mobility, its initial acceleration is $\nu_G$ times the corresponding component of
the final gradient.
\end{lrlemma}
\begin{proof}
At $V=0$, predicted probabilities are $u$, so the readout velocity is
$M$.  The gradient into $\phi$ has a factor $V$ and vanishes initially.
Differentiating the hidden gradient-flow equation at time zero leaves
$(D_\phi M)^\top M$, which is the gradient of $\LRnorm{M}^2/2$.
\end{proof}

For reference, if $P(\phi)=\mathbb E[(q_{s_{t-1}}-u)h_\phi^\top]$ is nonzero,
then the prior-orthogonalized rule feature
\begin{equation}
 R(\phi)=C(\phi)-\frac{\LRinner{C(\phi)}{P(\phi)}_F}
 {\LRnorm{P(\phi)}_F^2}P(\phi)
\end{equation}
satisfies
$\left.\frac{d}{dt}\LRinner{V(t)}{R(\phi(t))}_F\right|_{t=0}
=\LRnorm{R(\phi_0)}_F^2$.
This is an \emph{initial readout-alignment identity}, not a long-time
non-collapse theorem.

\begin{lrlemma}[Exact positional QK geometry]
\label{lem:lr-qk-geometry}
For the reduced positional projections $Q,K$ and
$S=QK^\top/\sqrt{d_h}$, let $G=\nabla_S\LRstep$.
When the positional embeddings are mutually orthogonal vectors
$\varepsilon e_i$ and the underlying positional QK matrix blocks share
mobility $\nu_Q$,
\begin{align}
 \dot Q&=-\frac{\nu_Q\varepsilon^2}{\sqrt{d_h}}GK,
 &\dot K&=-\frac{\nu_Q\varepsilon^2}{\sqrt{d_h}}G^\top Q,\label{eq:lr-qk-flow}\\
 \dot S&=-\frac{\nu_Q\varepsilon^2}{d_h}
 (QQ^\top G+GKK^\top),\label{eq:lr-qk-scoreflow}\\
 \frac{d}{dt}(Q^\top Q-K^\top K)&=0.\label{eq:lr-balanced-conservation}
\end{align}
These identities hold as readout and MLP parameters evolve, provided the
query/key dependence on the positional projections remains as stated.
\end{lrlemma}
\begin{proof}
Apply the chain rule to $QK^\top/\sqrt{d_h}$.  The one-hot positional
features supply the identical mobility factor $\nu_Q\varepsilon^2$ for
each projected row.  Substituting the two velocities in
$\frac{d}{dt}(Q^\top Q-K^\top K)$ cancels the four terms pairwise.
\end{proof}

In particular, if the conserved difference
$Q_0^\top Q_0-K_0^\top K_0$ has at least $n-1$ eigenvalues at least
$\mu>0$, the supervised-query Gram matrix $Q(t)Q(t)^\top$ has its
smallest eigenvalue at least $\mu$ as long as its number of rows is $n-1$.
This is a \emph{conditional} conditioning mechanism, not a conclusion for
arbitrarily initialized full Transformer QKV.

\subsection{A weakly planted feature without a planted route}
\label{app:weak-routing-family}

We now specify the symmetry-reduced family used by the remaining results.
Fix the query letter. Let $s=(-1)^{s_{t-1}}$ and let attention depend on positions, not the
random letter identities.  The full-prefix pooled, centered value is
\begin{equation}\label{eq:lr-centered-value}
 Z_t=\sum_{j\in\mathcal J}a_{tj}(b_j-p),\qquad
 a_{tj}=\frac{\exp(S_{i_tj})}{\sum_{k\le i_t}\exp(S_{i_tk})},
 \quad \mathcal J=\{j_1,\ldots,j_n\}.
\end{equation}
Non-letter positions have zero \emph{value}, not zero attention mass.
A fixed token/position-aware linear value direction realizes the centered contributions
$b_j-p$ for prompt letters and zero contributions on the remaining visible
sites.  The trainable value parameter is its amplitude $v$. For the fixed query, assign its token the usual match
value $1-p$ and subtract $1-p$ using its position embedding; assign count
and separator tokens value zero. Prompt-letter positions have zero positional
correction. Orthogonal token and position registers realize this exactly.
Allowing all
unrestricted value directions to vary is outside this invariant family.

A pair of ReLUs has preactivations
\begin{equation}\label{eq:lr-gates}
 r_+=\operatorname{ReLU}(as+bvZ_t+\beta),\qquad
 r_-=\operatorname{ReLU}(-as+bvZ_t+\beta).
\end{equation}
The binary logit (positive for state sign $+1$) is
\begin{equation}\label{eq:lr-factor-logit}
 f_t=cs+u\omega(r_+-r_-).
\end{equation}
Here $\beta$ is the \emph{MLP bias}, distinct from the attention
relative biases $\beta_\ell^{\mathrm{rel}}$ introduced below.
Every displayed previous state is available directly in the residual stream.

\begin{lrlemma}[Exact two-ReLU feature family]
\label{lem:lr-two-relu}
If
\begin{equation}\label{eq:lr-gate-condition}
 a>|\beta|+|bv|,
\end{equation}
then for all inputs with $|Z_t|\le1$,
\begin{equation}\label{eq:lr-feature-param}
 \begin{aligned}
 r_+-r_-&=s(d+wZ_t),& d&=a+\beta,& w&=bv,\\
 f_t&=s(C+gZ_t),& C&=c+u\omega d,& g&=u\omega w.
 \end{aligned}
\end{equation}
The feature family does not specify which prompt-letter position is current.
\end{lrlemma}
\begin{proof}
When $s=1$, only $r_+$ is active; when $s=-1$, only $r_-$ is active.
The inactive preactivation is strictly negative by
\eqref{eq:lr-gate-condition}.  Substitution gives the stated identity.
\end{proof}

\paragraph{Rule covariance and bounded perturbations.}
For a purely centered hidden coordinate $H_t=\rho sZ_t$,
Equations~\eqref{eq:lr-innovation-cov} and
\eqref{eq:lr-centered-value} give the exact identity
\begin{equation}\label{eq:lr-centered-seed-cov}
 \mathbb E[g_tH_t]=\rho\chi\bar A(e_1-e_0),\qquad
 \mathbb E[(q_{s_{t-1}}-u)H_t]=0.
\end{equation}
In particular the feature covariance is nonzero even when every
letter site has the same initial attention weight.  If a planted seed
has rule covariance at least $\kappa>0$ and the covariance map is
$L$-Lipschitz on a neighbourhood, then every perturbation of norm at
most $\kappa/(2L)$ preserves covariance at least $\kappa/2$.
This elementary robustness statement applies to perturbations within
a parameterization preserving the feature roles; it does not ensure
that arbitrary off-family trainable directions remain stable under
subsequent unrestricted gradient flow.

\paragraph{Initial positional routing acceleration (ideal centered channel).}
If the hidden representation contains the \emph{exact} independent
coordinate $H_t=\rho sZ_t$, all its other target-correlated coordinates
are independent of positional QK, and $V(0)=0$, then
$F(\phi)=\LRnorm{M(\phi)}_F^2/2$ satisfies, upon varying positional QK,
$F=F_0+\rho^2\chi^2\bar A^2$.  Thus under unit QK mobility,
\begin{equation}\label{eq:lr-centered-accel}
 \dot{\bar A}(0)=0,\qquad
 \ddot{\bar A}(0)=2\rho^2\chi^2\bar A(0)
 \LRnorm{\nabla_{Q,K}\bar A(0)}_F^2\ge0.
\end{equation}
Nondegenerate positional QK gives strict initial acceleration.
This statement is \emph{conditional on exact feature orthogonalization}:
the raw two-ReLU coordinate
$s(d+wZ_t)$ also contains prior information, which creates additional
terms.  It is therefore not a general acceleration theorem for arbitrary
planted Block 2 parameters.

The two hidden units may be trained as an \emph{explicitly tied Euclidean
submanifold}: opposite count weights, equal centered-letter weights and
biases, and opposite output weights.  The induced Euclidean metric assigns
factor $1/2$ to a gradient update of a scalar appearing twice with equal
magnitude.  This tied family is a model restriction; symmetry alone does
not prove its invariance for all unrestricted Transformer parameters.
The mobilities below define flow in the displayed effective factor coordinates.
If a matrix realization uses a fixed direction of nonunit norm, its norm
factor must be absorbed into the assigned coordinate mobility; the scalar
equations are not a claim about unnormalized full-matrix Euclidean flow.

Take $u(0)=c(0)=0$, $\omega_0>0$, $b_0<0$, $v_0>0$, and
$a_0>|\beta_0|+|b_0v_0|$.  Assume all planted magnitudes and gate margins
are inverse-polynomially bounded.  Initialize the positional attention
scores without a preferred correct position and set all trainable
relative-position biases to zero.  This is a \emph{weak computational seed},
not random discovery of the count--letter interaction.  Unlike the
fixed-sign-ReLU-bias initialization of Assumption~2 in the main paper,
the gates here deliberately respond differently to the displayed count.
Thus Theorem~3 does not automatically give a matched outcome plateau for
this precise planted initialization.

\subsection{Exact cross-entropy routing signs}
\label{app:weak-routing-signs}

Let $\sigma(x)=(1+e^{-x})^{-1}$, $y_t'=1-b_t$ and
\begin{equation}\label{eq:lr-logistic-risk}
 \LRstep=\sum_{t=2}^n\omega_t\,\mathbb E
 \left[\log(1+e^{C+gZ_t})-y_t'(C+gZ_t)\right].
\end{equation}
This is exactly the binary cross-entropy of
\eqref{eq:lr-feature-param} because multiplying the binary logit and
label by $s$ cancels the previous-state sign.
Let $A_t=a_{t,j_t}$ and $\bar A=\sum_t\omega_t A_t$.

\begin{lrlemma}[Innovation selects the current source]
\label{lem:lr-signs}
For position-only attention,
\begin{equation}\label{eq:lr-initial-gate-gradient}
 \left.\frac{\partial\LRstep}{\partial g}\right|_{g=0}
 =\chi\bar A>0
\end{equation}
for every $C$.  At any finite logits and $g<0$, define
$v_{tj}=\partial L_t/\partial a_{tj}$, where $L_t$ is the single-time
population risk without $\omega_t$.  Then
\begin{equation}\label{eq:lr-feature-source-signs}
 v_{t,j_t}<0,\qquad v_{tj}\ge0
 \quad(j\in\mathcal J,\;j\ne j_t),\qquad
 v_{tj}=0\quad(j\notin\mathcal J).
\end{equation}
If $|C|+|g|\le B$, then
\begin{equation}\label{eq:lr-source-bounds}
 v_{t,j_t}\le-2\chi|g|\sigma(-B),\qquad
 0\le v_{tj}\le\frac{\chi g^2a_{tj}}4\quad(j\ne j_t\text{ a letter}).
\end{equation}
\end{lrlemma}
\begin{proof}
Centering and independence give
$\mathbb E[(b_t-p)Z_t]=\chi A_t$, yielding
\eqref{eq:lr-initial-gate-gradient}.  For the correct letter,
conditioning on the other bits gives
\[
 v_{t,j_t}=g\chi\mathbb E
 \{1+\sigma(U+gA_t(1-p))-\sigma(U-gA_tp)\},
\]
where $U$ is the contribution of all other positions and $C$.
The braced term is strictly positive.  Under the logit bound it is at
least $2\sigma(-B)$.  For an incorrect letter $j$, the target is
independent of $b_j$, and conditioning on all other bits yields
\[
 v_{tj}=g\chi\mathbb E
 \{\sigma(U+ga_{tj}(1-p))-\sigma(U-ga_{tj}p)\}.
\]
For $g<0$ the two terms in braces are nonpositive; their magnitude is
at most $|g|a_{tj}/4$ by $\sigma'\le1/4$.  Non-letter values are zero.
\end{proof}

Put $G_{tj}=\omega_ta_{tj}(v_{tj}-\bar v_t)$, with
$\bar v_t=\sum_ka_{tk}v_{tk}$; this is the gradient with respect to a
pre-softmax score.  The following global comparison will be important.

\begin{lrlemma}[Correct-versus-incorrect score comparison]
\label{lem:lr-postrouting-gradient}
Whenever $g<0$ and $A_t\ge1/2$,
\begin{equation}\label{eq:lr-score-comparison}
 G_{t,j_t}<0,\qquad G_{tj}-G_{t,j_t}\ge0
 \quad\text{for every other visible site }j.
\end{equation}
Consequently, if the attention biases are shared by relative offset and
$i_t-j_t=n$, then for every $\ell\ne n$,
\begin{equation}\label{eq:lr-bias-gap-mono}
 \frac{d}{dt}
 (\beta_n^{\mathrm{rel}}-\beta_\ell^{\mathrm{rel}})\ge0
\end{equation}
under common positive relative-bias mobility.
\end{lrlemma}
\begin{proof}
The first inequality follows because the correct attention-weight
derivative is strictly smaller than all others, so it lies below their
weighted mean.  For an incorrect site with $a=a_{tj}\le A=A_t$ and
$v=v_{tj}\ge0$, $v_c=v_{t,j_t}<0$,
\[
 \frac{G_{tj}-G_{t,j_t}}{\omega_t}
 =av-Av_c+(A-a)\bar v_t
 \ge av+A(1-A+a)(-v_c)\ge0.
\]
We used $\bar v_t\ge Av_c$ and $A-a\ge0$.
An offset $\ell$ occurs at most once in any attention row.  The correct
offset $n$ appears in every supervised row.  Summing the comparisons
within rows where $\ell$ occurs, and using $G_{t,j_t}<0$ in the
remaining rows, gives
$\partial_{\beta_n^{\mathrm{rel}}}\LRstep
\le\partial_{\beta_\ell^{\mathrm{rel}}}\LRstep$.
\end{proof}

\subsection{Feature bootstrap with every parameter group active}
\label{app:weak-routing-bootstrap}

All training uses one \emph{fixed} system of positive mobilities:
$c,u$ have mobility one; tied MLP/value factors have mobility
$\nu_F>0$; positional QK have mobility $\nu_Q>0$; all relative biases
have common mobility $\nu_R>0$.
No parameter group is explicitly frozen or switched at a chosen time.
The factorized feature family and two-ReLU gate condition are nevertheless
restrictions on the parameterization.

Write $\alpha=u\omega$, $w=bv<0$, $d=a+\beta>0$,
$L_C=\partial\LRstep/\partial C$, $L_g=\partial\LRstep/\partial g$.
While the gate regime holds, the tied Euclidean gradient equations include
\begin{align}
 \dot c&=-L_C,
 &\dot u&=-\omega(dL_C+wL_g),\label{eq:lr-factor-u}\\
 \dot\omega&=-\tfrac{\nu_Fu}{2}(dL_C+wL_g),
 &\dot a=\dot\beta&=-\tfrac{\nu_F\alpha}{2}L_C,\\
 \dot b&=-\tfrac{\nu_F\alpha v}{2}L_g,
 &\dot v&=-\nu_F\alpha bL_g.
 \label{eq:lr-factor-bv}
\end{align}
In particular,
\begin{equation}\label{eq:lr-factor-invariants}
 u^2-\frac{2}{\nu_F}\omega^2=-\frac{2}{\nu_F}\omega_0^2,
 \qquad
 v^2-2b^2=v_0^2-2b_0^2,
 \qquad a-\beta=a_0-\beta_0.
\end{equation}
The first invariant gives $\omega^2\ge\omega_0^2$.  Under the gate
condition, the independent prior readout satisfies
$c(t)\le c_\star=\log((1-p)/p)$: at $c=c_\star$, with $u\omega\ge0$,
its derivative points inward.  At $u=0$ the velocity points into $u>0$,
and $w=bv$ retains its negative sign as long as the bounded gate regime
holds: on either boundary $b=0$ or $v=0$, one has $L_g=\chi\bar A>0$,
so the vector field points into $b<0,v>0$ (the simultaneous zero is excluded
by uniqueness from the nonzero seed). Thus $g=u\omega w<0$ for positive times.

\begin{lrlemma}[Quantitative weak-feature bootstrap, with explicit exits]
\label{lem:lr-feature-bootstrap}
Suppose on an interval prior to reaching the target magnitude that
\begin{equation}\label{eq:lr-amplification-region}
 \bar A\ge A_{\min}>0,\quad
 d\le D_F,\quad
 0<W\le|w|\le W_{\max},\quad
 \omega\ge\omega_{\min}>0,
\end{equation}
and the gate condition \eqref{eq:lr-gate-condition} holds.  Put
\begin{equation}\label{eq:lr-gamma-bootstrap}
 H=D_F+W_{\max},\quad
 \gamma=\chi A_{\min}\left(\frac W H\right)^2,
 \quad
 T_{\mathrm{boot}}=\frac4{3\omega_{\min}^2H^2}.
\end{equation}
Starting with $c=u=0$, the quantity $x=-g$ reaches $x\ge\gamma$
by $T_{\mathrm{boot}}$, unless one of the region conditions fails first.
No assumption $L_C\le0$ is required.
\end{lrlemma}
\begin{proof}
At $c\le c_\star$, let $\alpha=u\omega$.  Since
$0<d+wZ_t\le H$, logistic curvature implies
$L_C\le\alpha H/4$.  At $g=0$, $L_g=\chi\bar A$ and the
same curvature bound gives $L_g\ge\chi A_{\min}-x/4$.
When $x\le\gamma$,
\[
 dL_C+wL_g
 \le\frac{D_FH}{4W}x-W\chi A_{\min}
 +\frac{W_{\max}}4x
 \le-\frac34W\chi A_{\min}.
\]
The last step uses
$D_FH+WW_{\max}\le H^2$.
From \eqref{eq:lr-factor-u} and
$\omega^2\ge\omega_{\min}^2$,
$\dot\alpha\ge(3/4)\omega_{\min}^2W\chi A_{\min}$.
In this region $L_g>0$ and
$\dot w=-\nu_F\alpha(v^2+b^2/2)L_g\le0$.
Hence $\dot x\ge(3/4)\omega_{\min}^2W^2\chi A_{\min}$,
which gives the claimed hitting time.
\end{proof}

The groupwise loss budget is exact in the induced Euclidean metrics:
for the tied feature coordinates $(a,\beta,b,\omega,v)$, the squared
metric weights are $(2,2,2,2,1)$. For projected positional rows $(Q,K)$,
the squared metric weight is $\varepsilon^{-2}$.
\begin{equation}\label{eq:lr-mobility-budget}
 -\dot{\LRstep}=
 \LRnorm{\nabla_{c,u}\LRstep}^2
 +\nu_F\LRnorm{\nabla_F\LRstep}^2
 +\nu_Q\LRnorm{\nabla_{Q,K}\LRstep}^2
 +\nu_R\LRnorm{\nabla_{\beta^{\mathrm{rel}}}\LRstep}^2.
\end{equation}
Here group norms and gradients use those induced metrics. For each group
with mobility $\nu$, Cauchy--Schwarz implies
\begin{equation}\label{eq:lr-movement-budget}
 \LRnorm{\theta_G(T)-\theta_G(0)}
 \le\sqrt{\nu T\,[\LRstep(0)-\LRstep(T)]}
 \le\sqrt{\nu T\log2}.
\end{equation}
Thus sufficiently small but positive $\nu_F,\nu_Q$ preserve prescribed
strict feature-gate and base-score margins over any known finite horizon.
This justifies a controlled bootstrap; it does not establish unit-rate
stability in the unrestricted architecture.

\subsection{Relative-position learning and post-routing convergence}
\label{app:weak-routing-convergence}

Use positional scores
\begin{equation}\label{eq:lr-relative-scores}
 S_{i_tj}=B_{i_tj}+\beta_{i_t-j}^{\mathrm{rel}},
 \qquad
 B_{i_tj}=\frac{q_{i_t}^\top k_j}{\sqrt{d_h}},
 \qquad \beta_\ell^{\mathrm{rel}}(0)=0.
\end{equation}
The correct offset $n$ is not privileged at initialization.

\begin{lrlemma}[Correct-offset learning in the small-feature region]
\label{lem:lr-bias-learning}
Assume all base scores satisfy $|B_{i_tj}|\le D_A$,
all relative biases start at zero, and
$0<x=-g\le A_{\min}/4$. If $A_t\ge A_{\min}>0$,
then the correct score derivative is negative and all other visible-site
score derivatives are nonnegative:
\begin{equation}\label{eq:lr-smallscore-sign}
 G_{t,j_t}<0,\qquad G_{tj}\ge0\ (j\ne j_t).
\end{equation}
Consequently, for the shared relative biases,
\begin{equation}\label{eq:lr-beta-routes}
 \dot\beta_n^{\mathrm{rel}}>0,\qquad
 \dot\beta_\ell^{\mathrm{rel}}\le0\quad(\ell\ne n).
\end{equation}
They imply the uniform floor
\begin{equation}\label{eq:lr-floor}
 A_t\ge A_{\min}:=\frac{e^{-2D_A}}m,
 \qquad m:=\max_t\#\{j:j\le i_t\}\le3n.
\end{equation}
If the initial bootstrap reaches $x=\gamma$ and
$\gamma\le A_{\min}/4$, there is a polynomially bounded fixed bias
mobility $\nu_R$ that drives every $A_t$ to $2/3$ before $x$ leaves
$[\gamma/2,A_{\min}/4]$, provided the feature and base-score
neighbourhoods hold over that interval.
\end{lrlemma}
\begin{proof}
By Lemma~\ref{lem:lr-signs}, the correct value derivative is negative;
incorrect-letter derivatives are nonnegative.  The correct derivative has the logit-independent bound
\[
 -v_{t,j_t}\ge\chi x(1-xA_t/4)\ge\chi x(1-x/4),
\]
because $\sigma$ is $1/4$-Lipschitz. Hence
\[
 \bar v_t\le-\chi x A_t(1-x/4)+\chi x^2/4<0
\]
when $x\le A_{\min}/4$ and $A_t\ge A_{\min}$.  Hence the softmax derivatives have the stated signs.
The bias sign conclusion follows from the fixed offset
\eqref{eq:lr-constant-lag}.
With $\beta_n^{\mathrm{rel}}\ge0$, all other biases $\le0$,
and $|B|\le D_A$, softmax normalization yields
$A_t\ge e^{-2D_A}/m$.

For the time estimate, whenever some row has $A_t\le2/3$, the exact
correct-score derivative satisfies
$-G_{t,j_t}\ge c_0\omega_t\chi xA_{\min}$
with $c_0=1/4$ in this small-$x$ region: indeed
$-G_{t,j_t}\ge\omega_t A_t(1-A_t)(-v_{t,j_t})$.
Therefore
$\dot\beta_n^{\mathrm{rel}}
\ge c_0\nu_R\chi\gamma A_{\min}/[2(n-1)]$
as long as $x\ge\gamma/2$ and routing is incomplete.
All rows achieve $A_t\ge2/3$ when
$\beta_n^{\mathrm{rel}}\ge
2D_A+\log(2(m-1))$.
Thus the additional time is at most
\[
 T_{\mathrm{route}}
 \le \frac{2(n-1)[2D_A+\log(2(m-1))]}
 {c_0\nu_R\chi\gamma A_{\min}}.
\]
On a bounded feature neighbourhood, $|\dot x|\le M_x$ with $M_x$
independent of $\nu_R$ (attention probabilities stay in $[0,1]$).
Choose the fixed $\nu_R$ large enough that this time is less than
$\min\{1,\gamma/(4M_x),(A_{\min}/4-\gamma)/(2M_x)\}$.
For $p=1/4$, the choice $\gamma=\chi A_{\min}(W/H)^2$
automatically gives $\gamma\le3A_{\min}/16<A_{\min}/4$.
No learning-rate switch is required.
\end{proof}

\begin{lrlemma}[Post-routing invariance under bounded QK-score drift]
\label{lem:lr-post-routing}
Suppose $g<0$, $A_t(t_0)\ge2/3$ for all supervised $t$, and the
base-score contrasts change by at most $\log(5/4)$ after $t_0$:
\[
 |[B_{i_t,j_t}-B_{i_tj}](t)-[B_{i_t,j_t}-B_{i_tj}](t_0)|
 \le\log(5/4).
\]
As long as $g<0$ and the centered-feature family holds, all later
correct-position attention weights obey
\begin{equation}\label{eq:lr-post-floor}
 A_t(t)\ge\frac8{13}>\frac35.
\end{equation}
\end{lrlemma}
\begin{proof}
As long as all $A_t\ge1/2$, the bias comparisons of
Lemma~\ref{lem:lr-postrouting-gradient} imply that each relative-bias
gap $\beta_n^{\mathrm{rel}}-\beta_\ell^{\mathrm{rel}}$ is
nondecreasing.  The assumed bound on the other score contrasts says
the odds of each competitor relative to the correct site can increase
by at most a factor $5/4$.  Since the sum of competitor odds is at most
$1/2$ at time $t_0$, it stays at most $5/8$, yielding
$A_t\ge1/(1+5/8)=8/13$.  A first-exit argument closes the condition
$A_t\ge1/2$.
\end{proof}

\begin{lrlemma}[Uniform separability and convex-readout descent]
\label{lem:lr-readout-convergence}
Suppose $A_t\ge3/5$ at every supervised position.  Then
\begin{equation}\label{eq:lr-bit-separation}
 (1-2b_t)\left(\frac12-p-Z_t\right)\ge\frac1{10}
\end{equation}
for every prompt-bit realization.  Suppose in addition that, on a
finite interval after $t_0$, the trainable feature factors remain in a
neighbourhood that admits a \emph{fixed} readout comparator
$\widehat r=(\widehat c,\widehat u)$ satisfying
\begin{equation}\label{eq:lr-comparator}
 \LRstep(\widehat r;\phi(t))\le\eta/2
 \quad\text{for all }t\text{ in the interval}.
\end{equation}
If the two readout coefficients $(c,u)$ continue to train at unit
mobility and every other group has nonnegative mobility, then
\begin{equation}\label{eq:lr-descent}
 \LRstep(t_0+T)\le\frac\eta2+
 \frac{\LRnorm{r(t_0)-\widehat r}^2}{2T}.
\end{equation}
In particular, $T=\LRnorm{r(t_0)-\widehat r}^2/\eta$ suffices for
$\LRstep\le\eta$.
\end{lrlemma}
\begin{proof}
For $b_t=0$, $Z_t\le1-p-A_t$; for $b_t=1$,
$Z_t\ge A_t-p$.  Both cases imply
\eqref{eq:lr-bit-separation}.
For a reference seed $(d_0,w_0,\omega_0)$ with $w_0<0$, take
\[
 R=20\log(2/\eta),\qquad
 \widehat u=\frac{R}{\omega_0|w_0|},\qquad
 \widehat c=R\left(\frac12-p-\frac{d_0}{|w_0|}\right).
\]
At the reference feature, its effective logit is
$R(1/2-p-Z_t)$, which has signed margin at least $R/10$.
Require the feature neighbourhood to satisfy
\begin{equation}\label{eq:lr-comparator-stability}
 |\omega d-\omega_0d_0|+|\omega w-\omega_0w_0|
 \le\frac{\omega_0|w_0|}{20}.
\end{equation}
Since $|Z_t|\le1$, the logit perturbation is at most $R/20$.
The signed margin remains at least $R/20$, so the comparator loss is
at most $e^{-R/20}=\eta/2$. The same comparator is fixed before training,
independently of the routing hitting time.
Cross-entropy is convex in $r=(c,u)$ at fixed feature factors.  Thus
$\frac12\frac{d}{dt}\LRnorm{r-\widehat r}^2
\le\LRstep(\widehat r;\phi(t))-\LRstep(r;\phi(t))$.
The full joint flow makes the loss nonincreasing.  Integrating gives
\eqref{eq:lr-descent}.
\end{proof}

\begin{lrtheorem}[Finite-accuracy learned routing in a symmetry-reduced family]
\label{thm:lr-conditional-bridge}
Fix $n\ge2$, $p=1/4$, $0<\eta<1/2$ and $0<\varepsilon\le1$.
Use the explicitly tied family \eqref{eq:lr-gates}--\eqref{eq:lr-feature-param}
with $c_0=u_0=0$, $b_0<0$, $v_0,\omega_0>0$, and
$a_0>|\beta_0|+|b_0v_0|$. The fixed query and centered value direction
are as in \cref{app:weak-routing-family}. QK consists only of trainable
positional projections of orthogonal position embeddings; every visible
causal site participates in softmax. All relative biases start at zero.

Assume a complexity bound $P\ge2$ bounds $n$, the widths,
$\varepsilon^{-1}$, all initial factor magnitudes, and the Frobenius norms
of the projected QK matrices. Also require
\[
 \min\{|b_0|,v_0,\omega_0,a_0-|\beta_0|-|b_0v_0|\}\ge P^{-1}.
\]
The initial base scores must have moderate spread: for some $D_0\ge0$,
\[
 |B_{i_tj}(0)|\le D_0,\qquad e^{2D_0}\le P.
\]
No correct relative offset is privileged by the relative biases. The base
scores may be bounded perturbations of equal scores; the assertion does
not require an initially favorable correct-site score.

There exist fixed positive mobilities $\nu_F,\nu_Q$ and $\nu_R$, with
$\nu_F^{-1},\nu_Q^{-1},\nu_R$ bounded by polynomials in $P$ and $1/\eta$,
such that the uninterrupted joint population gradient flow reaches
\begin{equation}\label{eq:lr-finite-theorem}
 \LRstep(T)\le\eta
\end{equation}
for $T\le\operatorname{poly}(P,1/\eta)$.
Before the post-routing convergence interval it reaches $A_t\ge2/3$
for every supervised $t$, and preserves $A_t\ge8/13$ throughout that
interval. The mobilities are chosen once at initialization, as below.

The claim concerns a \emph{target-dependent finite-accuracy schedule of
constant group mobilities}, not asymptotic loss convergence under one
fixed mobility choice independent of $\eta$, and not the unrestricted
Transformer parameterization.
\end{lrtheorem}
\begin{proof}[Proof with mobilities fixed before training]
Choose feature and QK balls of radii $r_F,r_Q>0$ around their initial
coordinates. Their inverse radii can be polynomially bounded in $P$.
For the feature ball require the gate condition, $b<0$, $v,\omega>0$,
\eqref{eq:lr-comparator-stability}, and uniform bounds
\[
 0<W\le|w|\le W_{\max},\quad d\le D_F,\quad
 \omega_{\min}\le\omega\le\Omega.
\]
The quantities $W^{-1},\omega_{\min}^{-1},W_{\max},D_F,\Omega$ are
polynomial in $P$. For the QK ball require $|B|\le D_A=D_0+1$ and
\[
 |(B_{i_t,j_t}-B_{i_tj})-(B_{i_t,j_t}(0)-B_{i_tj}(0))|
 \le\tfrac12\log(5/4).
\]
Polynomial maps on bounded seed balls are Lipschitz with polynomial
constants; the assumed inverse-polynomial gate and factor margins give
these radii. Both QK points at any two times in the ball therefore have
contrast differences at most $\log(5/4)$.

Set $A_{\min}=e^{-2D_A}/m$, $H=D_F+W_{\max}$,
$\gamma=\chi A_{\min}(W/H)^2$ and $T_b=T_{\mathrm{boot}}$ from
\eqref{eq:lr-gamma-bootstrap}. These quantities and their relevant
inverses are polynomial in $P$, and $\gamma\le3A_{\min}/16$.
Set $T_{\mathrm{pre}}=T_b+1$. Until that horizon, inside the feature
ball, the unit-mobility readout equations give
\[
 |\dot c|\le1,\qquad |\dot u|\le\Omega H,
 \qquad \|(c,u)\|\le T_{\mathrm{pre}}\sqrt{1+\Omega^2H^2}.
\]
These bounds use only $|\sigma-y'|\le1$ and $|Z_t|\le1$ and do not
depend on bias mobility. With $\nu_F\le1$, the factor equations also
give a polynomial bound $M_x\ge1$ on $|\dot x|$ in this interval,
uniformly over every attention distribution. Choose
\[
 \tau=\min\{1,\gamma/(4M_x),(A_{\min}/4-\gamma)/(2M_x)\},
 \quad K_R=2D_A+\log(2(m-1)),
\]
\[
 \nu_R=\max\left\{1,
 \frac{2(n-1)K_R}{c_0\chi\gamma A_{\min}\tau}\right\},
 \qquad c_0=1/4.
\]
This fixes the relative-bias mobility in advance and makes the routing
bound at most $\tau$. It is polynomially bounded.

Use the fixed reference comparator from the readout lemma and define
\[
 D=T_{\mathrm{pre}}\sqrt{1+\Omega^2H^2}+\|\widehat r\|,
 \qquad T_{\mathrm{post}}=D^2/\eta,\qquad
 T_*=T_{\mathrm{pre}}+T_{\mathrm{post}}.
\]
Now fix the other mobilities by
\[
 \nu_F=\min\{1,r_F^2/(4T_*\log2)\},\qquad
 \nu_Q=\min\{1,r_Q^2/(4T_*\log2)\}.
\]
Their inverses are polynomially bounded. For the tied feature metric,
all squared weights are at least one; for projected QK the weight is
$\varepsilon^{-2}\ge1$. Thus the movement budget bounds ordinary
coordinate movement by $r_F/2$ and $r_Q/2$ up to $T_*$.
A first-exit argument therefore keeps both groups strictly inside their
chosen balls for the whole argument. Finite-time existence follows from
the smooth loss and its gradient-flow energy budget: all mobilities are
finite positive constants, so finite-time parameter escape is excluded.

Inside these balls, $c\le c_\star$ and $u\ge0$ persist by the inward
boundary signs stated above; $\omega>0$ and $w<0$ give $g<0$ for positive
times. In the small-feature region the bias signs imply the attention
floor by a first-exit argument from zero biases. The bootstrap lemma
reaches $x=\gamma$ by $T_b$ unless all routing is already complete;
even in that case the same bootstrap bound applies. Starting at this
hitting time, within at most $\tau$ the bias lemma makes every
$A_t\ge2/3$ (if this already holds, the extra time is zero). Let $t_0$
be the end of this interval; then $t_0\le T_{\mathrm{pre}}$.
The derivative bound keeps $x$ inside $[\gamma/2,A_{\min}/4]$ during
this routing interval. The bounded base-score contrasts and the
post-routing lemma preserve $A_t\ge8/13>3/5$ afterwards.
The feature ball preserves the fixed comparator for the entire remaining
interval. Its initial distance at $t_0$ is at most $D$, so the readout
lemma gives $\LRstep(t_0+T_{\mathrm{post}})\le\eta$.
All mobilities were selected from initial uniform ball bounds; none
requires a later switch or a posterior choice of $t_0$.
\end{proof}

\subsection{Limitations and relationship to the main comparison}
\label{app:weak-routing-limitations}

\paragraph{What is new.}
The statistical signal for the correct relative offset is derived from
independent innovations rather than inserted in the attention mask or bias
initialization.  All visible causal sites compete in the softmax.  The
tied two-ReLU coefficients, value amplitude, positional QK, linear readout,
and relative-position biases train \emph{simultaneously}, although at
different fixed mobilities.  The constant-offset identity
\eqref{eq:lr-constant-lag} is a genuine architectural inductive bias.

\paragraph{What is not proved.}
The reduced family preselects a state--letter feature direction and
centers prompt-letter values.  It excludes unrestricted content-QKV
updates and arbitrary perturbations that leave this tied family.
The result also permits a large relative-bias mobility and very small
positive feature/QK mobilities, and both may depend on the target error
$\eta$.  This is therefore weaker than a theorem for ordinary
unit-mobility, randomly initialized Block 2.  In addition, the
count-dependent ReLU activity differs from the fixed-sign near-zero
bias regime in the main paper's Assumption~2.  The original outcome
plateau (Theorem~3) cannot be transferred to this planted seed without
an additional argument.  Nor does the finite-accuracy guarantee imply
$\LRstep(t)\to0$ for a single mobility choice independent of $\eta$.

\paragraph{Scope of the audited parameterization.}
For interpreting this appendix extension, the tied-subfamily
symmetry, parameter scaling relative to Assumption~2, and value-role restriction must be distinguished
from unrestricted updates in the original computational graph. The quantitative construction above
uses the stated reduced-coordinate metrics, moderate base-score spread,
and bounded seeds; it does not transfer to arbitrary parameter metrics.  In particular, the
centering of the fixed query-letter token must be implemented by the
value projection rather than ignored, the ordinary token-content QK
directions must be either controlled or explicitly excluded, and any
extension to Theorem~3 must separately verify its activation-pattern
hypotheses.  An earlier inference that \emph{nonzero rule covariance by
itself} forces correct routing is false; count-prior interference can
make the initial positional routing acceleration negative.

\paragraph{Suggested interpretation.}
This result is a complementary \emph{conditional routing bridge} next
to the original Theorem~4, not a replacement of its rigorous
positive-margin/global-loss convergence statement.  The next genuinely
stronger theorem would remove the tied feature restriction while
keeping a planted feature, and then establish feature stability,
relative-position learning, and polynomial-time low loss for all
trainable Block 2 directions, with a matched outcome comparison.

\paragraph{Numerical audit.} Exact Bernoulli enumeration and automatic differentiation in
\nolinkurl{handcoded/theory_checks/learned_routing_extension_checks.py}
check the source gradients, small-feature and post-routing score signs, tied-factor
invariants, positional QK metric, uniform separability, and the fixed comparator.
These checks support the displayed local identities and inequalities; they are
not a formal verification of the global theorem or a test of unrestricted training.
 in the manuscript
% appendix. Requires amsmath, amssymb, amsthm (already common in NeurIPS papers).
%
% STATUS: Audited reduced-model extension; quantitative hypotheses are explicit.
% It is an appendix addition and does not alter either part of Theorem 4.
% The end-to-end conclusion is explicitly CONDITIONAL on a symmetry-preserving
% planted-feature parameterization and quantitative parameter-group mobilities.
% It is NOT a theorem for unrestricted random Transformer initialization.
%
% Suggested main-text replacement/addition for Section 4.2 (copy separately):
% "Theorem 4 isolates transition learning with supplied useful directions and,
% in the clean binary extension, learns distractor suppression. A complementary
% full-prefix analysis (Appendix [reference]) considers a weakly seeded
% count--letter feature with no prescribed current-letter position. In a
% symmetry-reduced two-ReLU model, the innovation covariance selects the
% correct relative offset, and fixed-mobility joint gradient flow can learn
% this bias and reach low step loss in polynomial time. This extension still
% relies on tied feature directions, a small parameter-group learning rate,
% and an accelerated trainable relative bias; it does not prove convergence
% for unrestricted Block 2 training or by itself extend Theorem 3's matched
% outcome plateau to the planted activation-pattern initialization."
% ============================================================================

\providecommand{\LRstep}{L_{\mathrm{step}}}
\providecommand{\LRout}{L_{\mathrm{out}}}
\providecommand{\LRsig}{\sigma}
\providecommand{\LRnorm}[1]{\left\lVert #1\right\rVert}
\providecommand{\LRinner}[2]{\left\langle #1,#2\right\rangle}
\theoremstyle{plain}
\newtheorem{lrlemma}[lemma]{Lemma}
\newtheorem{lrproposition}[proposition]{Proposition}
\newtheorem{lrtheorem}[theorem]{Theorem}
\newtheorem{lrremark}[remark]{Remark}
\crefname{lrlemma}{Lemma}{Lemmas}
\crefname{lrproposition}{Proposition}{Propositions}
\crefname{lrtheorem}{Theorem}{Theorems}
\crefname{lrremark}{Remark}{Remarks}

\section{Weakly seeded full-prefix routing with joint feature learning}
\label{app:weak-seeded-full-prefix-routing}

\paragraph{Relation to Theorem 4.}
Theorem~4 of the main paper provides a polynomial-time positive-margin
convergence theorem with prescribed useful attention directions, and a clean
binary variant that learns soft penalties on designated zero-key/value
distractors.  This section studies a different intermediate regime: the model
is supplied a weak \emph{count--letter feature}, but no current-letter position
or preferred relative offset.  Attention remains causal and full-prefix;
relative-position biases are all initially zero.  The purpose is to isolate
which parts of routing and transition learning follow from population
cross-entropy and which require additional architectural restrictions.
Throughout, the terms ``full-prefix'' and ``joint training'' do not mean
unrestricted Transformer parameters: a tied, symmetry-reduced value/MLP
family is maintained as specified below.

\subsection{Data, representation, and the exact initial identities}
\label{app:weak-routing-initial}

Consider binary counting $K=2$, with independent
$b_j\sim\operatorname{Bernoulli}(p)$, $0<p<1/2$, and uniform $s_0$.
The application of interest is $p=1/4$.  As in the main paper, the input
contains the prompt letters and the displayed running counts, and the
prediction at $i_t$ for $s_t$ sees $s_{t-1}$; the associated prompt-letter
site is $j_t$ and
\begin{equation}\label{eq:lr-constant-lag}
 i_t-j_t=n,\qquad t=2,\ldots,n.
\end{equation}
The attention distribution is normalized over the \emph{whole visible causal
prefix}, including query, separator, and displayed-state sites.  We write
$\omega_t=1/(n-1)$ and $\chi=p(1-p)$.

For the actual one-head Block 2 with hidden parameters $\phi$ and readout
$V$, let $h_t=h_\phi(x_t)$, $z_t=Vh_t$, and
\begin{equation}
 q_s=(1-p)e_s+pe_{s\oplus1},\qquad
 g_t=e_{s_t}-q_{s_{t-1}}=(b_t-p)(e_{s_{t-1}\oplus1}-e_{s_{t-1}}).
\end{equation}
The state $s_{t-1}$ is independent of the entire prompt-bit vector: this
follows from uniformity of $s_0$.  Consequently,
\begin{equation}\label{eq:lr-innovation-cov}
 \mathbb E[(b_t-p)(b_j-p)]=\chi\mathbf 1\{j=t\}.
\end{equation}
For fixed $p=1/4$, the right side is a positive constant, unlike the
outcome-correlation factor $\lambda^n$.

\begin{lrlemma}[Exact target decomposition and readout rule credit]
\label{lem:lr-decomposition}
For population step cross-entropy and arbitrary logits,
\begin{equation}\label{eq:lr-kl-decomp}
 \LRstep=H(p)+\mathbb E\operatorname{KL}
 (q_{s_{t-1}}\Vert\operatorname{softmax}(z_t))
 -\Phi_{\mathrm{rule}},\qquad
 \Phi_{\mathrm{rule}}=\mathbb E\LRinner{g_t}{z_t}.
\end{equation}
Writing $C(\phi)=\mathbb E[g_t h_t^\top]$, one has
$\nabla_V\Phi_{\mathrm{rule}}=C(\phi)$.  At an exact state-only predictor
$\operatorname{softmax}(z_t)=q_{s_{t-1}}$, one has
$\nabla_V\LRstep=-C(\phi)$.
\end{lrlemma}
\begin{proof}
The conditional expected label is $q_{s_{t-1}}$ and the centered
innovation is $g_t$.  Expanding cross-entropy as log-sum-exp minus the
target logit gives \eqref{eq:lr-kl-decomp}.  The readout derivatives follow
from $z_t=Vh_t$ and the usual softmax cross-entropy gradient.
\end{proof}

\begin{lrlemma}[Exact zero-readout bootstrap]
\label{lem:lr-zero-bootstrap}
Let $u=(1/2,1/2)$ and
$M(\phi)=\mathbb E[(e_{s_t}-u)h_\phi^\top]$.
Under \emph{unit-mobility Euclidean gradient flow} of $(\phi,V)$ with
$V(0)=0$, at any initial parameter point with a locally continuously differentiable hidden map,
\begin{equation}\label{eq:lr-zero-boot}
 \dot V(0)=M(\phi_0),\qquad \dot\phi(0)=0,\qquad
 \ddot\phi(0)=\nabla_\phi\frac12\LRnorm{M(\phi_0)}_F^2.
\end{equation}
For a group of hidden parameters trained with constant mobility $\nu_G$,
with unit readout mobility, its initial acceleration is $\nu_G$ times the corresponding component of
the final gradient.
\end{lrlemma}
\begin{proof}
At $V=0$, predicted probabilities are $u$, so the readout velocity is
$M$.  The gradient into $\phi$ has a factor $V$ and vanishes initially.
Differentiating the hidden gradient-flow equation at time zero leaves
$(D_\phi M)^\top M$, which is the gradient of $\LRnorm{M}^2/2$.
\end{proof}

For reference, if $P(\phi)=\mathbb E[(q_{s_{t-1}}-u)h_\phi^\top]$ is nonzero,
then the prior-orthogonalized rule feature
\begin{equation}
 R(\phi)=C(\phi)-\frac{\LRinner{C(\phi)}{P(\phi)}_F}
 {\LRnorm{P(\phi)}_F^2}P(\phi)
\end{equation}
satisfies
$\left.\frac{d}{dt}\LRinner{V(t)}{R(\phi(t))}_F\right|_{t=0}
=\LRnorm{R(\phi_0)}_F^2$.
This is an \emph{initial readout-alignment identity}, not a long-time
non-collapse theorem.

\begin{lrlemma}[Exact positional QK geometry]
\label{lem:lr-qk-geometry}
For the reduced positional projections $Q,K$ and
$S=QK^\top/\sqrt{d_h}$, let $G=\nabla_S\LRstep$.
When the positional embeddings are mutually orthogonal vectors
$\varepsilon e_i$ and the underlying positional QK matrix blocks share
mobility $\nu_Q$,
\begin{align}
 \dot Q&=-\frac{\nu_Q\varepsilon^2}{\sqrt{d_h}}GK,
 &\dot K&=-\frac{\nu_Q\varepsilon^2}{\sqrt{d_h}}G^\top Q,\label{eq:lr-qk-flow}\\
 \dot S&=-\frac{\nu_Q\varepsilon^2}{d_h}
 (QQ^\top G+GKK^\top),\label{eq:lr-qk-scoreflow}\\
 \frac{d}{dt}(Q^\top Q-K^\top K)&=0.\label{eq:lr-balanced-conservation}
\end{align}
These identities hold as readout and MLP parameters evolve, provided the
query/key dependence on the positional projections remains as stated.
\end{lrlemma}
\begin{proof}
Apply the chain rule to $QK^\top/\sqrt{d_h}$.  The one-hot positional
features supply the identical mobility factor $\nu_Q\varepsilon^2$ for
each projected row.  Substituting the two velocities in
$\frac{d}{dt}(Q^\top Q-K^\top K)$ cancels the four terms pairwise.
\end{proof}

In particular, if the conserved difference
$Q_0^\top Q_0-K_0^\top K_0$ has at least $n-1$ eigenvalues at least
$\mu>0$, the supervised-query Gram matrix $Q(t)Q(t)^\top$ has its
smallest eigenvalue at least $\mu$ as long as its number of rows is $n-1$.
This is a \emph{conditional} conditioning mechanism, not a conclusion for
arbitrarily initialized full Transformer QKV.

\subsection{A weakly planted feature without a planted route}
\label{app:weak-routing-family}

We now specify the symmetry-reduced family used by the remaining results.
Fix the query letter. Let $s=(-1)^{s_{t-1}}$ and let attention depend on positions, not the
random letter identities.  The full-prefix pooled, centered value is
\begin{equation}\label{eq:lr-centered-value}
 Z_t=\sum_{j\in\mathcal J}a_{tj}(b_j-p),\qquad
 a_{tj}=\frac{\exp(S_{i_tj})}{\sum_{k\le i_t}\exp(S_{i_tk})},
 \quad \mathcal J=\{j_1,\ldots,j_n\}.
\end{equation}
Non-letter positions have zero \emph{value}, not zero attention mass.
A fixed token/position-aware linear value direction realizes the centered contributions
$b_j-p$ for prompt letters and zero contributions on the remaining visible
sites.  The trainable value parameter is its amplitude $v$. For the fixed query, assign its token the usual match
value $1-p$ and subtract $1-p$ using its position embedding; assign count
and separator tokens value zero. Prompt-letter positions have zero positional
correction. Orthogonal token and position registers realize this exactly.
Allowing all
unrestricted value directions to vary is outside this invariant family.

A pair of ReLUs has preactivations
\begin{equation}\label{eq:lr-gates}
 r_+=\operatorname{ReLU}(as+bvZ_t+\beta),\qquad
 r_-=\operatorname{ReLU}(-as+bvZ_t+\beta).
\end{equation}
The binary logit (positive for state sign $+1$) is
\begin{equation}\label{eq:lr-factor-logit}
 f_t=cs+u\omega(r_+-r_-).
\end{equation}
Here $\beta$ is the \emph{MLP bias}, distinct from the attention
relative biases $\beta_\ell^{\mathrm{rel}}$ introduced below.
Every displayed previous state is available directly in the residual stream.

\begin{lrlemma}[Exact two-ReLU feature family]
\label{lem:lr-two-relu}
If
\begin{equation}\label{eq:lr-gate-condition}
 a>|\beta|+|bv|,
\end{equation}
then for all inputs with $|Z_t|\le1$,
\begin{equation}\label{eq:lr-feature-param}
 \begin{aligned}
 r_+-r_-&=s(d+wZ_t),& d&=a+\beta,& w&=bv,\\
 f_t&=s(C+gZ_t),& C&=c+u\omega d,& g&=u\omega w.
 \end{aligned}
\end{equation}
The feature family does not specify which prompt-letter position is current.
\end{lrlemma}
\begin{proof}
When $s=1$, only $r_+$ is active; when $s=-1$, only $r_-$ is active.
The inactive preactivation is strictly negative by
\eqref{eq:lr-gate-condition}.  Substitution gives the stated identity.
\end{proof}

\paragraph{Rule covariance and bounded perturbations.}
For a purely centered hidden coordinate $H_t=\rho sZ_t$,
Equations~\eqref{eq:lr-innovation-cov} and
\eqref{eq:lr-centered-value} give the exact identity
\begin{equation}\label{eq:lr-centered-seed-cov}
 \mathbb E[g_tH_t]=\rho\chi\bar A(e_1-e_0),\qquad
 \mathbb E[(q_{s_{t-1}}-u)H_t]=0.
\end{equation}
In particular the feature covariance is nonzero even when every
letter site has the same initial attention weight.  If a planted seed
has rule covariance at least $\kappa>0$ and the covariance map is
$L$-Lipschitz on a neighbourhood, then every perturbation of norm at
most $\kappa/(2L)$ preserves covariance at least $\kappa/2$.
This elementary robustness statement applies to perturbations within
a parameterization preserving the feature roles; it does not ensure
that arbitrary off-family trainable directions remain stable under
subsequent unrestricted gradient flow.

\paragraph{Initial positional routing acceleration (ideal centered channel).}
If the hidden representation contains the \emph{exact} independent
coordinate $H_t=\rho sZ_t$, all its other target-correlated coordinates
are independent of positional QK, and $V(0)=0$, then
$F(\phi)=\LRnorm{M(\phi)}_F^2/2$ satisfies, upon varying positional QK,
$F=F_0+\rho^2\chi^2\bar A^2$.  Thus under unit QK mobility,
\begin{equation}\label{eq:lr-centered-accel}
 \dot{\bar A}(0)=0,\qquad
 \ddot{\bar A}(0)=2\rho^2\chi^2\bar A(0)
 \LRnorm{\nabla_{Q,K}\bar A(0)}_F^2\ge0.
\end{equation}
Nondegenerate positional QK gives strict initial acceleration.
This statement is \emph{conditional on exact feature orthogonalization}:
the raw two-ReLU coordinate
$s(d+wZ_t)$ also contains prior information, which creates additional
terms.  It is therefore not a general acceleration theorem for arbitrary
planted Block 2 parameters.

The two hidden units may be trained as an \emph{explicitly tied Euclidean
submanifold}: opposite count weights, equal centered-letter weights and
biases, and opposite output weights.  The induced Euclidean metric assigns
factor $1/2$ to a gradient update of a scalar appearing twice with equal
magnitude.  This tied family is a model restriction; symmetry alone does
not prove its invariance for all unrestricted Transformer parameters.
The mobilities below define flow in the displayed effective factor coordinates.
If a matrix realization uses a fixed direction of nonunit norm, its norm
factor must be absorbed into the assigned coordinate mobility; the scalar
equations are not a claim about unnormalized full-matrix Euclidean flow.

Take $u(0)=c(0)=0$, $\omega_0>0$, $b_0<0$, $v_0>0$, and
$a_0>|\beta_0|+|b_0v_0|$.  Assume all planted magnitudes and gate margins
are inverse-polynomially bounded.  Initialize the positional attention
scores without a preferred correct position and set all trainable
relative-position biases to zero.  This is a \emph{weak computational seed},
not random discovery of the count--letter interaction.  Unlike the
fixed-sign-ReLU-bias initialization of Assumption~2 in the main paper,
the gates here deliberately respond differently to the displayed count.
Thus Theorem~3 does not automatically give a matched outcome plateau for
this precise planted initialization.

\subsection{Exact cross-entropy routing signs}
\label{app:weak-routing-signs}

Let $\sigma(x)=(1+e^{-x})^{-1}$, $y_t'=1-b_t$ and
\begin{equation}\label{eq:lr-logistic-risk}
 \LRstep=\sum_{t=2}^n\omega_t\,\mathbb E
 \left[\log(1+e^{C+gZ_t})-y_t'(C+gZ_t)\right].
\end{equation}
This is exactly the binary cross-entropy of
\eqref{eq:lr-feature-param} because multiplying the binary logit and
label by $s$ cancels the previous-state sign.
Let $A_t=a_{t,j_t}$ and $\bar A=\sum_t\omega_t A_t$.

\begin{lrlemma}[Innovation selects the current source]
\label{lem:lr-signs}
For position-only attention,
\begin{equation}\label{eq:lr-initial-gate-gradient}
 \left.\frac{\partial\LRstep}{\partial g}\right|_{g=0}
 =\chi\bar A>0
\end{equation}
for every $C$.  At any finite logits and $g<0$, define
$v_{tj}=\partial L_t/\partial a_{tj}$, where $L_t$ is the single-time
population risk without $\omega_t$.  Then
\begin{equation}\label{eq:lr-feature-source-signs}
 v_{t,j_t}<0,\qquad v_{tj}\ge0
 \quad(j\in\mathcal J,\;j\ne j_t),\qquad
 v_{tj}=0\quad(j\notin\mathcal J).
\end{equation}
If $|C|+|g|\le B$, then
\begin{equation}\label{eq:lr-source-bounds}
 v_{t,j_t}\le-2\chi|g|\sigma(-B),\qquad
 0\le v_{tj}\le\frac{\chi g^2a_{tj}}4\quad(j\ne j_t\text{ a letter}).
\end{equation}
\end{lrlemma}
\begin{proof}
Centering and independence give
$\mathbb E[(b_t-p)Z_t]=\chi A_t$, yielding
\eqref{eq:lr-initial-gate-gradient}.  For the correct letter,
conditioning on the other bits gives
\[
 v_{t,j_t}=g\chi\mathbb E
 \{1+\sigma(U+gA_t(1-p))-\sigma(U-gA_tp)\},
\]
where $U$ is the contribution of all other positions and $C$.
The braced term is strictly positive.  Under the logit bound it is at
least $2\sigma(-B)$.  For an incorrect letter $j$, the target is
independent of $b_j$, and conditioning on all other bits yields
\[
 v_{tj}=g\chi\mathbb E
 \{\sigma(U+ga_{tj}(1-p))-\sigma(U-ga_{tj}p)\}.
\]
For $g<0$ the two terms in braces are nonpositive; their magnitude is
at most $|g|a_{tj}/4$ by $\sigma'\le1/4$.  Non-letter values are zero.
\end{proof}

Put $G_{tj}=\omega_ta_{tj}(v_{tj}-\bar v_t)$, with
$\bar v_t=\sum_ka_{tk}v_{tk}$; this is the gradient with respect to a
pre-softmax score.  The following global comparison will be important.

\begin{lrlemma}[Correct-versus-incorrect score comparison]
\label{lem:lr-postrouting-gradient}
Whenever $g<0$ and $A_t\ge1/2$,
\begin{equation}\label{eq:lr-score-comparison}
 G_{t,j_t}<0,\qquad G_{tj}-G_{t,j_t}\ge0
 \quad\text{for every other visible site }j.
\end{equation}
Consequently, if the attention biases are shared by relative offset and
$i_t-j_t=n$, then for every $\ell\ne n$,
\begin{equation}\label{eq:lr-bias-gap-mono}
 \frac{d}{dt}
 (\beta_n^{\mathrm{rel}}-\beta_\ell^{\mathrm{rel}})\ge0
\end{equation}
under common positive relative-bias mobility.
\end{lrlemma}
\begin{proof}
The first inequality follows because the correct attention-weight
derivative is strictly smaller than all others, so it lies below their
weighted mean.  For an incorrect site with $a=a_{tj}\le A=A_t$ and
$v=v_{tj}\ge0$, $v_c=v_{t,j_t}<0$,
\[
 \frac{G_{tj}-G_{t,j_t}}{\omega_t}
 =av-Av_c+(A-a)\bar v_t
 \ge av+A(1-A+a)(-v_c)\ge0.
\]
We used $\bar v_t\ge Av_c$ and $A-a\ge0$.
An offset $\ell$ occurs at most once in any attention row.  The correct
offset $n$ appears in every supervised row.  Summing the comparisons
within rows where $\ell$ occurs, and using $G_{t,j_t}<0$ in the
remaining rows, gives
$\partial_{\beta_n^{\mathrm{rel}}}\LRstep
\le\partial_{\beta_\ell^{\mathrm{rel}}}\LRstep$.
\end{proof}

\subsection{Feature bootstrap with every parameter group active}
\label{app:weak-routing-bootstrap}

All training uses one \emph{fixed} system of positive mobilities:
$c,u$ have mobility one; tied MLP/value factors have mobility
$\nu_F>0$; positional QK have mobility $\nu_Q>0$; all relative biases
have common mobility $\nu_R>0$.
No parameter group is explicitly frozen or switched at a chosen time.
The factorized feature family and two-ReLU gate condition are nevertheless
restrictions on the parameterization.

Write $\alpha=u\omega$, $w=bv<0$, $d=a+\beta>0$,
$L_C=\partial\LRstep/\partial C$, $L_g=\partial\LRstep/\partial g$.
While the gate regime holds, the tied Euclidean gradient equations include
\begin{align}
 \dot c&=-L_C,
 &\dot u&=-\omega(dL_C+wL_g),\label{eq:lr-factor-u}\\
 \dot\omega&=-\tfrac{\nu_Fu}{2}(dL_C+wL_g),
 &\dot a=\dot\beta&=-\tfrac{\nu_F\alpha}{2}L_C,\\
 \dot b&=-\tfrac{\nu_F\alpha v}{2}L_g,
 &\dot v&=-\nu_F\alpha bL_g.
 \label{eq:lr-factor-bv}
\end{align}
In particular,
\begin{equation}\label{eq:lr-factor-invariants}
 u^2-\frac{2}{\nu_F}\omega^2=-\frac{2}{\nu_F}\omega_0^2,
 \qquad
 v^2-2b^2=v_0^2-2b_0^2,
 \qquad a-\beta=a_0-\beta_0.
\end{equation}
The first invariant gives $\omega^2\ge\omega_0^2$.  Under the gate
condition, the independent prior readout satisfies
$c(t)\le c_\star=\log((1-p)/p)$: at $c=c_\star$, with $u\omega\ge0$,
its derivative points inward.  At $u=0$ the velocity points into $u>0$,
and $w=bv$ retains its negative sign as long as the bounded gate regime
holds: on either boundary $b=0$ or $v=0$, one has $L_g=\chi\bar A>0$,
so the vector field points into $b<0,v>0$ (the simultaneous zero is excluded
by uniqueness from the nonzero seed). Thus $g=u\omega w<0$ for positive times.

\begin{lrlemma}[Quantitative weak-feature bootstrap, with explicit exits]
\label{lem:lr-feature-bootstrap}
Suppose on an interval prior to reaching the target magnitude that
\begin{equation}\label{eq:lr-amplification-region}
 \bar A\ge A_{\min}>0,\quad
 d\le D_F,\quad
 0<W\le|w|\le W_{\max},\quad
 \omega\ge\omega_{\min}>0,
\end{equation}
and the gate condition \eqref{eq:lr-gate-condition} holds.  Put
\begin{equation}\label{eq:lr-gamma-bootstrap}
 H=D_F+W_{\max},\quad
 \gamma=\chi A_{\min}\left(\frac W H\right)^2,
 \quad
 T_{\mathrm{boot}}=\frac4{3\omega_{\min}^2H^2}.
\end{equation}
Starting with $c=u=0$, the quantity $x=-g$ reaches $x\ge\gamma$
by $T_{\mathrm{boot}}$, unless one of the region conditions fails first.
No assumption $L_C\le0$ is required.
\end{lrlemma}
\begin{proof}
At $c\le c_\star$, let $\alpha=u\omega$.  Since
$0<d+wZ_t\le H$, logistic curvature implies
$L_C\le\alpha H/4$.  At $g=0$, $L_g=\chi\bar A$ and the
same curvature bound gives $L_g\ge\chi A_{\min}-x/4$.
When $x\le\gamma$,
\[
 dL_C+wL_g
 \le\frac{D_FH}{4W}x-W\chi A_{\min}
 +\frac{W_{\max}}4x
 \le-\frac34W\chi A_{\min}.
\]
The last step uses
$D_FH+WW_{\max}\le H^2$.
From \eqref{eq:lr-factor-u} and
$\omega^2\ge\omega_{\min}^2$,
$\dot\alpha\ge(3/4)\omega_{\min}^2W\chi A_{\min}$.
In this region $L_g>0$ and
$\dot w=-\nu_F\alpha(v^2+b^2/2)L_g\le0$.
Hence $\dot x\ge(3/4)\omega_{\min}^2W^2\chi A_{\min}$,
which gives the claimed hitting time.
\end{proof}

The groupwise loss budget is exact in the induced Euclidean metrics:
for the tied feature coordinates $(a,\beta,b,\omega,v)$, the squared
metric weights are $(2,2,2,2,1)$. For projected positional rows $(Q,K)$,
the squared metric weight is $\varepsilon^{-2}$.
\begin{equation}\label{eq:lr-mobility-budget}
 -\dot{\LRstep}=
 \LRnorm{\nabla_{c,u}\LRstep}^2
 +\nu_F\LRnorm{\nabla_F\LRstep}^2
 +\nu_Q\LRnorm{\nabla_{Q,K}\LRstep}^2
 +\nu_R\LRnorm{\nabla_{\beta^{\mathrm{rel}}}\LRstep}^2.
\end{equation}
Here group norms and gradients use those induced metrics. For each group
with mobility $\nu$, Cauchy--Schwarz implies
\begin{equation}\label{eq:lr-movement-budget}
 \LRnorm{\theta_G(T)-\theta_G(0)}
 \le\sqrt{\nu T\,[\LRstep(0)-\LRstep(T)]}
 \le\sqrt{\nu T\log2}.
\end{equation}
Thus sufficiently small but positive $\nu_F,\nu_Q$ preserve prescribed
strict feature-gate and base-score margins over any known finite horizon.
This justifies a controlled bootstrap; it does not establish unit-rate
stability in the unrestricted architecture.

\subsection{Relative-position learning and post-routing convergence}
\label{app:weak-routing-convergence}

Use positional scores
\begin{equation}\label{eq:lr-relative-scores}
 S_{i_tj}=B_{i_tj}+\beta_{i_t-j}^{\mathrm{rel}},
 \qquad
 B_{i_tj}=\frac{q_{i_t}^\top k_j}{\sqrt{d_h}},
 \qquad \beta_\ell^{\mathrm{rel}}(0)=0.
\end{equation}
The correct offset $n$ is not privileged at initialization.

\begin{lrlemma}[Correct-offset learning in the small-feature region]
\label{lem:lr-bias-learning}
Assume all base scores satisfy $|B_{i_tj}|\le D_A$,
all relative biases start at zero, and
$0<x=-g\le A_{\min}/4$. If $A_t\ge A_{\min}>0$,
then the correct score derivative is negative and all other visible-site
score derivatives are nonnegative:
\begin{equation}\label{eq:lr-smallscore-sign}
 G_{t,j_t}<0,\qquad G_{tj}\ge0\ (j\ne j_t).
\end{equation}
Consequently, for the shared relative biases,
\begin{equation}\label{eq:lr-beta-routes}
 \dot\beta_n^{\mathrm{rel}}>0,\qquad
 \dot\beta_\ell^{\mathrm{rel}}\le0\quad(\ell\ne n).
\end{equation}
They imply the uniform floor
\begin{equation}\label{eq:lr-floor}
 A_t\ge A_{\min}:=\frac{e^{-2D_A}}m,
 \qquad m:=\max_t\#\{j:j\le i_t\}\le3n.
\end{equation}
If the initial bootstrap reaches $x=\gamma$ and
$\gamma\le A_{\min}/4$, there is a polynomially bounded fixed bias
mobility $\nu_R$ that drives every $A_t$ to $2/3$ before $x$ leaves
$[\gamma/2,A_{\min}/4]$, provided the feature and base-score
neighbourhoods hold over that interval.
\end{lrlemma}
\begin{proof}
By Lemma~\ref{lem:lr-signs}, the correct value derivative is negative;
incorrect-letter derivatives are nonnegative.  The correct derivative has the logit-independent bound
\[
 -v_{t,j_t}\ge\chi x(1-xA_t/4)\ge\chi x(1-x/4),
\]
because $\sigma$ is $1/4$-Lipschitz. Hence
\[
 \bar v_t\le-\chi x A_t(1-x/4)+\chi x^2/4<0
\]
when $x\le A_{\min}/4$ and $A_t\ge A_{\min}$.  Hence the softmax derivatives have the stated signs.
The bias sign conclusion follows from the fixed offset
\eqref{eq:lr-constant-lag}.
With $\beta_n^{\mathrm{rel}}\ge0$, all other biases $\le0$,
and $|B|\le D_A$, softmax normalization yields
$A_t\ge e^{-2D_A}/m$.

For the time estimate, whenever some row has $A_t\le2/3$, the exact
correct-score derivative satisfies
$-G_{t,j_t}\ge c_0\omega_t\chi xA_{\min}$
with $c_0=1/4$ in this small-$x$ region: indeed
$-G_{t,j_t}\ge\omega_t A_t(1-A_t)(-v_{t,j_t})$.
Therefore
$\dot\beta_n^{\mathrm{rel}}
\ge c_0\nu_R\chi\gamma A_{\min}/[2(n-1)]$
as long as $x\ge\gamma/2$ and routing is incomplete.
All rows achieve $A_t\ge2/3$ when
$\beta_n^{\mathrm{rel}}\ge
2D_A+\log(2(m-1))$.
Thus the additional time is at most
\[
 T_{\mathrm{route}}
 \le \frac{2(n-1)[2D_A+\log(2(m-1))]}
 {c_0\nu_R\chi\gamma A_{\min}}.
\]
On a bounded feature neighbourhood, $|\dot x|\le M_x$ with $M_x$
independent of $\nu_R$ (attention probabilities stay in $[0,1]$).
Choose the fixed $\nu_R$ large enough that this time is less than
$\min\{1,\gamma/(4M_x),(A_{\min}/4-\gamma)/(2M_x)\}$.
For $p=1/4$, the choice $\gamma=\chi A_{\min}(W/H)^2$
automatically gives $\gamma\le3A_{\min}/16<A_{\min}/4$.
No learning-rate switch is required.
\end{proof}

\begin{lrlemma}[Post-routing invariance under bounded QK-score drift]
\label{lem:lr-post-routing}
Suppose $g<0$, $A_t(t_0)\ge2/3$ for all supervised $t$, and the
base-score contrasts change by at most $\log(5/4)$ after $t_0$:
\[
 |[B_{i_t,j_t}-B_{i_tj}](t)-[B_{i_t,j_t}-B_{i_tj}](t_0)|
 \le\log(5/4).
\]
As long as $g<0$ and the centered-feature family holds, all later
correct-position attention weights obey
\begin{equation}\label{eq:lr-post-floor}
 A_t(t)\ge\frac8{13}>\frac35.
\end{equation}
\end{lrlemma}
\begin{proof}
As long as all $A_t\ge1/2$, the bias comparisons of
Lemma~\ref{lem:lr-postrouting-gradient} imply that each relative-bias
gap $\beta_n^{\mathrm{rel}}-\beta_\ell^{\mathrm{rel}}$ is
nondecreasing.  The assumed bound on the other score contrasts says
the odds of each competitor relative to the correct site can increase
by at most a factor $5/4$.  Since the sum of competitor odds is at most
$1/2$ at time $t_0$, it stays at most $5/8$, yielding
$A_t\ge1/(1+5/8)=8/13$.  A first-exit argument closes the condition
$A_t\ge1/2$.
\end{proof}

\begin{lrlemma}[Uniform separability and convex-readout descent]
\label{lem:lr-readout-convergence}
Suppose $A_t\ge3/5$ at every supervised position.  Then
\begin{equation}\label{eq:lr-bit-separation}
 (1-2b_t)\left(\frac12-p-Z_t\right)\ge\frac1{10}
\end{equation}
for every prompt-bit realization.  Suppose in addition that, on a
finite interval after $t_0$, the trainable feature factors remain in a
neighbourhood that admits a \emph{fixed} readout comparator
$\widehat r=(\widehat c,\widehat u)$ satisfying
\begin{equation}\label{eq:lr-comparator}
 \LRstep(\widehat r;\phi(t))\le\eta/2
 \quad\text{for all }t\text{ in the interval}.
\end{equation}
If the two readout coefficients $(c,u)$ continue to train at unit
mobility and every other group has nonnegative mobility, then
\begin{equation}\label{eq:lr-descent}
 \LRstep(t_0+T)\le\frac\eta2+
 \frac{\LRnorm{r(t_0)-\widehat r}^2}{2T}.
\end{equation}
In particular, $T=\LRnorm{r(t_0)-\widehat r}^2/\eta$ suffices for
$\LRstep\le\eta$.
\end{lrlemma}
\begin{proof}
For $b_t=0$, $Z_t\le1-p-A_t$; for $b_t=1$,
$Z_t\ge A_t-p$.  Both cases imply
\eqref{eq:lr-bit-separation}.
For a reference seed $(d_0,w_0,\omega_0)$ with $w_0<0$, take
\[
 R=20\log(2/\eta),\qquad
 \widehat u=\frac{R}{\omega_0|w_0|},\qquad
 \widehat c=R\left(\frac12-p-\frac{d_0}{|w_0|}\right).
\]
At the reference feature, its effective logit is
$R(1/2-p-Z_t)$, which has signed margin at least $R/10$.
Require the feature neighbourhood to satisfy
\begin{equation}\label{eq:lr-comparator-stability}
 |\omega d-\omega_0d_0|+|\omega w-\omega_0w_0|
 \le\frac{\omega_0|w_0|}{20}.
\end{equation}
Since $|Z_t|\le1$, the logit perturbation is at most $R/20$.
The signed margin remains at least $R/20$, so the comparator loss is
at most $e^{-R/20}=\eta/2$. The same comparator is fixed before training,
independently of the routing hitting time.
Cross-entropy is convex in $r=(c,u)$ at fixed feature factors.  Thus
$\frac12\frac{d}{dt}\LRnorm{r-\widehat r}^2
\le\LRstep(\widehat r;\phi(t))-\LRstep(r;\phi(t))$.
The full joint flow makes the loss nonincreasing.  Integrating gives
\eqref{eq:lr-descent}.
\end{proof}

\begin{lrtheorem}[Finite-accuracy learned routing in a symmetry-reduced family]
\label{thm:lr-conditional-bridge}
Fix $n\ge2$, $p=1/4$, $0<\eta<1/2$ and $0<\varepsilon\le1$.
Use the explicitly tied family \eqref{eq:lr-gates}--\eqref{eq:lr-feature-param}
with $c_0=u_0=0$, $b_0<0$, $v_0,\omega_0>0$, and
$a_0>|\beta_0|+|b_0v_0|$. The fixed query and centered value direction
are as in \cref{app:weak-routing-family}. QK consists only of trainable
positional projections of orthogonal position embeddings; every visible
causal site participates in softmax. All relative biases start at zero.

Assume a complexity bound $P\ge2$ bounds $n$, the widths,
$\varepsilon^{-1}$, all initial factor magnitudes, and the Frobenius norms
of the projected QK matrices. Also require
\[
 \min\{|b_0|,v_0,\omega_0,a_0-|\beta_0|-|b_0v_0|\}\ge P^{-1}.
\]
The initial base scores must have moderate spread: for some $D_0\ge0$,
\[
 |B_{i_tj}(0)|\le D_0,\qquad e^{2D_0}\le P.
\]
No correct relative offset is privileged by the relative biases. The base
scores may be bounded perturbations of equal scores; the assertion does
not require an initially favorable correct-site score.

There exist fixed positive mobilities $\nu_F,\nu_Q$ and $\nu_R$, with
$\nu_F^{-1},\nu_Q^{-1},\nu_R$ bounded by polynomials in $P$ and $1/\eta$,
such that the uninterrupted joint population gradient flow reaches
\begin{equation}\label{eq:lr-finite-theorem}
 \LRstep(T)\le\eta
\end{equation}
for $T\le\operatorname{poly}(P,1/\eta)$.
Before the post-routing convergence interval it reaches $A_t\ge2/3$
for every supervised $t$, and preserves $A_t\ge8/13$ throughout that
interval. The mobilities are chosen once at initialization, as below.

The claim concerns a \emph{target-dependent finite-accuracy schedule of
constant group mobilities}, not asymptotic loss convergence under one
fixed mobility choice independent of $\eta$, and not the unrestricted
Transformer parameterization.
\end{lrtheorem}
\begin{proof}[Proof with mobilities fixed before training]
Choose feature and QK balls of radii $r_F,r_Q>0$ around their initial
coordinates. Their inverse radii can be polynomially bounded in $P$.
For the feature ball require the gate condition, $b<0$, $v,\omega>0$,
\eqref{eq:lr-comparator-stability}, and uniform bounds
\[
 0<W\le|w|\le W_{\max},\quad d\le D_F,\quad
 \omega_{\min}\le\omega\le\Omega.
\]
The quantities $W^{-1},\omega_{\min}^{-1},W_{\max},D_F,\Omega$ are
polynomial in $P$. For the QK ball require $|B|\le D_A=D_0+1$ and
\[
 |(B_{i_t,j_t}-B_{i_tj})-(B_{i_t,j_t}(0)-B_{i_tj}(0))|
 \le\tfrac12\log(5/4).
\]
Polynomial maps on bounded seed balls are Lipschitz with polynomial
constants; the assumed inverse-polynomial gate and factor margins give
these radii. Both QK points at any two times in the ball therefore have
contrast differences at most $\log(5/4)$.

Set $A_{\min}=e^{-2D_A}/m$, $H=D_F+W_{\max}$,
$\gamma=\chi A_{\min}(W/H)^2$ and $T_b=T_{\mathrm{boot}}$ from
\eqref{eq:lr-gamma-bootstrap}. These quantities and their relevant
inverses are polynomial in $P$, and $\gamma\le3A_{\min}/16$.
Set $T_{\mathrm{pre}}=T_b+1$. Until that horizon, inside the feature
ball, the unit-mobility readout equations give
\[
 |\dot c|\le1,\qquad |\dot u|\le\Omega H,
 \qquad \|(c,u)\|\le T_{\mathrm{pre}}\sqrt{1+\Omega^2H^2}.
\]
These bounds use only $|\sigma-y'|\le1$ and $|Z_t|\le1$ and do not
depend on bias mobility. With $\nu_F\le1$, the factor equations also
give a polynomial bound $M_x\ge1$ on $|\dot x|$ in this interval,
uniformly over every attention distribution. Choose
\[
 \tau=\min\{1,\gamma/(4M_x),(A_{\min}/4-\gamma)/(2M_x)\},
 \quad K_R=2D_A+\log(2(m-1)),
\]
\[
 \nu_R=\max\left\{1,
 \frac{2(n-1)K_R}{c_0\chi\gamma A_{\min}\tau}\right\},
 \qquad c_0=1/4.
\]
This fixes the relative-bias mobility in advance and makes the routing
bound at most $\tau$. It is polynomially bounded.

Use the fixed reference comparator from the readout lemma and define
\[
 D=T_{\mathrm{pre}}\sqrt{1+\Omega^2H^2}+\|\widehat r\|,
 \qquad T_{\mathrm{post}}=D^2/\eta,\qquad
 T_*=T_{\mathrm{pre}}+T_{\mathrm{post}}.
\]
Now fix the other mobilities by
\[
 \nu_F=\min\{1,r_F^2/(4T_*\log2)\},\qquad
 \nu_Q=\min\{1,r_Q^2/(4T_*\log2)\}.
\]
Their inverses are polynomially bounded. For the tied feature metric,
all squared weights are at least one; for projected QK the weight is
$\varepsilon^{-2}\ge1$. Thus the movement budget bounds ordinary
coordinate movement by $r_F/2$ and $r_Q/2$ up to $T_*$.
A first-exit argument therefore keeps both groups strictly inside their
chosen balls for the whole argument. Finite-time existence follows from
the smooth loss and its gradient-flow energy budget: all mobilities are
finite positive constants, so finite-time parameter escape is excluded.

Inside these balls, $c\le c_\star$ and $u\ge0$ persist by the inward
boundary signs stated above; $\omega>0$ and $w<0$ give $g<0$ for positive
times. In the small-feature region the bias signs imply the attention
floor by a first-exit argument from zero biases. The bootstrap lemma
reaches $x=\gamma$ by $T_b$ unless all routing is already complete;
even in that case the same bootstrap bound applies. Starting at this
hitting time, within at most $\tau$ the bias lemma makes every
$A_t\ge2/3$ (if this already holds, the extra time is zero). Let $t_0$
be the end of this interval; then $t_0\le T_{\mathrm{pre}}$.
The derivative bound keeps $x$ inside $[\gamma/2,A_{\min}/4]$ during
this routing interval. The bounded base-score contrasts and the
post-routing lemma preserve $A_t\ge8/13>3/5$ afterwards.
The feature ball preserves the fixed comparator for the entire remaining
interval. Its initial distance at $t_0$ is at most $D$, so the readout
lemma gives $\LRstep(t_0+T_{\mathrm{post}})\le\eta$.
All mobilities were selected from initial uniform ball bounds; none
requires a later switch or a posterior choice of $t_0$.
\end{proof}

\subsection{Limitations and relationship to the main comparison}
\label{app:weak-routing-limitations}

\paragraph{What is new.}
The statistical signal for the correct relative offset is derived from
independent innovations rather than inserted in the attention mask or bias
initialization.  All visible causal sites compete in the softmax.  The
tied two-ReLU coefficients, value amplitude, positional QK, linear readout,
and relative-position biases train \emph{simultaneously}, although at
different fixed mobilities.  The constant-offset identity
\eqref{eq:lr-constant-lag} is a genuine architectural inductive bias.

\paragraph{What is not proved.}
The reduced family preselects a state--letter feature direction and
centers prompt-letter values.  It excludes unrestricted content-QKV
updates and arbitrary perturbations that leave this tied family.
The result also permits a large relative-bias mobility and very small
positive feature/QK mobilities, and both may depend on the target error
$\eta$.  This is therefore weaker than a theorem for ordinary
unit-mobility, randomly initialized Block 2.  In addition, the
count-dependent ReLU activity differs from the fixed-sign near-zero
bias regime in the main paper's Assumption~2.  The original outcome
plateau (Theorem~3) cannot be transferred to this planted seed without
an additional argument.  Nor does the finite-accuracy guarantee imply
$\LRstep(t)\to0$ for a single mobility choice independent of $\eta$.

\paragraph{Scope of the audited parameterization.}
For interpreting this appendix extension, the tied-subfamily
symmetry, parameter scaling relative to Assumption~2, and value-role restriction must be distinguished
from unrestricted updates in the original computational graph. The quantitative construction above
uses the stated reduced-coordinate metrics, moderate base-score spread,
and bounded seeds; it does not transfer to arbitrary parameter metrics.  In particular, the
centering of the fixed query-letter token must be implemented by the
value projection rather than ignored, the ordinary token-content QK
directions must be either controlled or explicitly excluded, and any
extension to Theorem~3 must separately verify its activation-pattern
hypotheses.  An earlier inference that \emph{nonzero rule covariance by
itself} forces correct routing is false; count-prior interference can
make the initial positional routing acceleration negative.

\paragraph{Suggested interpretation.}
This result is a complementary \emph{conditional routing bridge} next
to the original Theorem~4, not a replacement of its rigorous
positive-margin/global-loss convergence statement.  The next genuinely
stronger theorem would remove the tied feature restriction while
keeping a planted feature, and then establish feature stability,
relative-position learning, and polynomial-time low loss for all
trainable Block 2 directions, with a matched outcome comparison.

\paragraph{Numerical audit.} Exact Bernoulli enumeration and automatic differentiation in
\nolinkurl{handcoded/theory_checks/learned_routing_extension_checks.py}
check the source gradients, small-feature and post-routing score signs, tied-factor
invariants, positional QK metric, uniform separability, and the fixed comparator.
These checks support the displayed local identities and inequalities; they are
not a formal verification of the global theorem or a test of unrestricted training.
 in the manuscript
% appendix. Requires amsmath, amssymb, amsthm (already common in NeurIPS papers).
%
% STATUS: Audited reduced-model extension; quantitative hypotheses are explicit.
% It is an appendix addition and does not alter either part of Theorem 4.
% The end-to-end conclusion is explicitly CONDITIONAL on a symmetry-preserving
% planted-feature parameterization and quantitative parameter-group mobilities.
% It is NOT a theorem for unrestricted random Transformer initialization.
%
% Suggested main-text replacement/addition for Section 4.2 (copy separately):
% "Theorem 4 isolates transition learning with supplied useful directions and,
% in the clean binary extension, learns distractor suppression. A complementary
% full-prefix analysis (Appendix [reference]) considers a weakly seeded
% count--letter feature with no prescribed current-letter position. In a
% symmetry-reduced two-ReLU model, the innovation covariance selects the
% correct relative offset, and fixed-mobility joint gradient flow can learn
% this bias and reach low step loss in polynomial time. This extension still
% relies on tied feature directions, a small parameter-group learning rate,
% and an accelerated trainable relative bias; it does not prove convergence
% for unrestricted Block 2 training or by itself extend Theorem 3's matched
% outcome plateau to the planted activation-pattern initialization."
% ============================================================================

\providecommand{\LRstep}{L_{\mathrm{step}}}
\providecommand{\LRout}{L_{\mathrm{out}}}
\providecommand{\LRsig}{\sigma}
\providecommand{\LRnorm}[1]{\left\lVert #1\right\rVert}
\providecommand{\LRinner}[2]{\left\langle #1,#2\right\rangle}
\theoremstyle{plain}
\newtheorem{lrlemma}[lemma]{Lemma}
\newtheorem{lrproposition}[proposition]{Proposition}
\newtheorem{lrtheorem}[theorem]{Theorem}
\newtheorem{lrremark}[remark]{Remark}
\crefname{lrlemma}{Lemma}{Lemmas}
\crefname{lrproposition}{Proposition}{Propositions}
\crefname{lrtheorem}{Theorem}{Theorems}
\crefname{lrremark}{Remark}{Remarks}

\section{Weakly seeded full-prefix routing with joint feature learning}
\label{app:weak-seeded-full-prefix-routing}

\paragraph{Relation to Theorem 4.}
Theorem~4 of the main paper provides a polynomial-time positive-margin
convergence theorem with prescribed useful attention directions, and a clean
binary variant that learns soft penalties on designated zero-key/value
distractors.  This section studies a different intermediate regime: the model
is supplied a weak \emph{count--letter feature}, but no current-letter position
or preferred relative offset.  Attention remains causal and full-prefix;
relative-position biases are all initially zero.  The purpose is to isolate
which parts of routing and transition learning follow from population
cross-entropy and which require additional architectural restrictions.
Throughout, the terms ``full-prefix'' and ``joint training'' do not mean
unrestricted Transformer parameters: a tied, symmetry-reduced value/MLP
family is maintained as specified below.

\subsection{Data, representation, and the exact initial identities}
\label{app:weak-routing-initial}

Consider binary counting $K=2$, with independent
$b_j\sim\operatorname{Bernoulli}(p)$, $0<p<1/2$, and uniform $s_0$.
The application of interest is $p=1/4$.  As in the main paper, the input
contains the prompt letters and the displayed running counts, and the
prediction at $i_t$ for $s_t$ sees $s_{t-1}$; the associated prompt-letter
site is $j_t$ and
\begin{equation}\label{eq:lr-constant-lag}
 i_t-j_t=n,\qquad t=2,\ldots,n.
\end{equation}
The attention distribution is normalized over the \emph{whole visible causal
prefix}, including query, separator, and displayed-state sites.  We write
$\omega_t=1/(n-1)$ and $\chi=p(1-p)$.

For the actual one-head Block 2 with hidden parameters $\phi$ and readout
$V$, let $h_t=h_\phi(x_t)$, $z_t=Vh_t$, and
\begin{equation}
 q_s=(1-p)e_s+pe_{s\oplus1},\qquad
 g_t=e_{s_t}-q_{s_{t-1}}=(b_t-p)(e_{s_{t-1}\oplus1}-e_{s_{t-1}}).
\end{equation}
The state $s_{t-1}$ is independent of the entire prompt-bit vector: this
follows from uniformity of $s_0$.  Consequently,
\begin{equation}\label{eq:lr-innovation-cov}
 \mathbb E[(b_t-p)(b_j-p)]=\chi\mathbf 1\{j=t\}.
\end{equation}
For fixed $p=1/4$, the right side is a positive constant, unlike the
outcome-correlation factor $\lambda^n$.

\begin{lrlemma}[Exact target decomposition and readout rule credit]
\label{lem:lr-decomposition}
For population step cross-entropy and arbitrary logits,
\begin{equation}\label{eq:lr-kl-decomp}
 \LRstep=H(p)+\mathbb E\operatorname{KL}
 (q_{s_{t-1}}\Vert\operatorname{softmax}(z_t))
 -\Phi_{\mathrm{rule}},\qquad
 \Phi_{\mathrm{rule}}=\mathbb E\LRinner{g_t}{z_t}.
\end{equation}
Writing $C(\phi)=\mathbb E[g_t h_t^\top]$, one has
$\nabla_V\Phi_{\mathrm{rule}}=C(\phi)$.  At an exact state-only predictor
$\operatorname{softmax}(z_t)=q_{s_{t-1}}$, one has
$\nabla_V\LRstep=-C(\phi)$.
\end{lrlemma}
\begin{proof}
The conditional expected label is $q_{s_{t-1}}$ and the centered
innovation is $g_t$.  Expanding cross-entropy as log-sum-exp minus the
target logit gives \eqref{eq:lr-kl-decomp}.  The readout derivatives follow
from $z_t=Vh_t$ and the usual softmax cross-entropy gradient.
\end{proof}

\begin{lrlemma}[Exact zero-readout bootstrap]
\label{lem:lr-zero-bootstrap}
Let $u=(1/2,1/2)$ and
$M(\phi)=\mathbb E[(e_{s_t}-u)h_\phi^\top]$.
Under \emph{unit-mobility Euclidean gradient flow} of $(\phi,V)$ with
$V(0)=0$, at any initial parameter point with a locally continuously differentiable hidden map,
\begin{equation}\label{eq:lr-zero-boot}
 \dot V(0)=M(\phi_0),\qquad \dot\phi(0)=0,\qquad
 \ddot\phi(0)=\nabla_\phi\frac12\LRnorm{M(\phi_0)}_F^2.
\end{equation}
For a group of hidden parameters trained with constant mobility $\nu_G$,
with unit readout mobility, its initial acceleration is $\nu_G$ times the corresponding component of
the final gradient.
\end{lrlemma}
\begin{proof}
At $V=0$, predicted probabilities are $u$, so the readout velocity is
$M$.  The gradient into $\phi$ has a factor $V$ and vanishes initially.
Differentiating the hidden gradient-flow equation at time zero leaves
$(D_\phi M)^\top M$, which is the gradient of $\LRnorm{M}^2/2$.
\end{proof}

For reference, if $P(\phi)=\mathbb E[(q_{s_{t-1}}-u)h_\phi^\top]$ is nonzero,
then the prior-orthogonalized rule feature
\begin{equation}
 R(\phi)=C(\phi)-\frac{\LRinner{C(\phi)}{P(\phi)}_F}
 {\LRnorm{P(\phi)}_F^2}P(\phi)
\end{equation}
satisfies
$\left.\frac{d}{dt}\LRinner{V(t)}{R(\phi(t))}_F\right|_{t=0}
=\LRnorm{R(\phi_0)}_F^2$.
This is an \emph{initial readout-alignment identity}, not a long-time
non-collapse theorem.

\begin{lrlemma}[Exact positional QK geometry]
\label{lem:lr-qk-geometry}
For the reduced positional projections $Q,K$ and
$S=QK^\top/\sqrt{d_h}$, let $G=\nabla_S\LRstep$.
When the positional embeddings are mutually orthogonal vectors
$\varepsilon e_i$ and the underlying positional QK matrix blocks share
mobility $\nu_Q$,
\begin{align}
 \dot Q&=-\frac{\nu_Q\varepsilon^2}{\sqrt{d_h}}GK,
 &\dot K&=-\frac{\nu_Q\varepsilon^2}{\sqrt{d_h}}G^\top Q,\label{eq:lr-qk-flow}\\
 \dot S&=-\frac{\nu_Q\varepsilon^2}{d_h}
 (QQ^\top G+GKK^\top),\label{eq:lr-qk-scoreflow}\\
 \frac{d}{dt}(Q^\top Q-K^\top K)&=0.\label{eq:lr-balanced-conservation}
\end{align}
These identities hold as readout and MLP parameters evolve, provided the
query/key dependence on the positional projections remains as stated.
\end{lrlemma}
\begin{proof}
Apply the chain rule to $QK^\top/\sqrt{d_h}$.  The one-hot positional
features supply the identical mobility factor $\nu_Q\varepsilon^2$ for
each projected row.  Substituting the two velocities in
$\frac{d}{dt}(Q^\top Q-K^\top K)$ cancels the four terms pairwise.
\end{proof}

In particular, if the conserved difference
$Q_0^\top Q_0-K_0^\top K_0$ has at least $n-1$ eigenvalues at least
$\mu>0$, the supervised-query Gram matrix $Q(t)Q(t)^\top$ has its
smallest eigenvalue at least $\mu$ as long as its number of rows is $n-1$.
This is a \emph{conditional} conditioning mechanism, not a conclusion for
arbitrarily initialized full Transformer QKV.

\subsection{A weakly planted feature without a planted route}
\label{app:weak-routing-family}

We now specify the symmetry-reduced family used by the remaining results.
Fix the query letter. Let $s=(-1)^{s_{t-1}}$ and let attention depend on positions, not the
random letter identities.  The full-prefix pooled, centered value is
\begin{equation}\label{eq:lr-centered-value}
 Z_t=\sum_{j\in\mathcal J}a_{tj}(b_j-p),\qquad
 a_{tj}=\frac{\exp(S_{i_tj})}{\sum_{k\le i_t}\exp(S_{i_tk})},
 \quad \mathcal J=\{j_1,\ldots,j_n\}.
\end{equation}
Non-letter positions have zero \emph{value}, not zero attention mass.
A fixed token/position-aware linear value direction realizes the centered contributions
$b_j-p$ for prompt letters and zero contributions on the remaining visible
sites.  The trainable value parameter is its amplitude $v$. For the fixed query, assign its token the usual match
value $1-p$ and subtract $1-p$ using its position embedding; assign count
and separator tokens value zero. Prompt-letter positions have zero positional
correction. Orthogonal token and position registers realize this exactly.
Allowing all
unrestricted value directions to vary is outside this invariant family.

A pair of ReLUs has preactivations
\begin{equation}\label{eq:lr-gates}
 r_+=\operatorname{ReLU}(as+bvZ_t+\beta),\qquad
 r_-=\operatorname{ReLU}(-as+bvZ_t+\beta).
\end{equation}
The binary logit (positive for state sign $+1$) is
\begin{equation}\label{eq:lr-factor-logit}
 f_t=cs+u\omega(r_+-r_-).
\end{equation}
Here $\beta$ is the \emph{MLP bias}, distinct from the attention
relative biases $\beta_\ell^{\mathrm{rel}}$ introduced below.
Every displayed previous state is available directly in the residual stream.

\begin{lrlemma}[Exact two-ReLU feature family]
\label{lem:lr-two-relu}
If
\begin{equation}\label{eq:lr-gate-condition}
 a>|\beta|+|bv|,
\end{equation}
then for all inputs with $|Z_t|\le1$,
\begin{equation}\label{eq:lr-feature-param}
 \begin{aligned}
 r_+-r_-&=s(d+wZ_t),& d&=a+\beta,& w&=bv,\\
 f_t&=s(C+gZ_t),& C&=c+u\omega d,& g&=u\omega w.
 \end{aligned}
\end{equation}
The feature family does not specify which prompt-letter position is current.
\end{lrlemma}
\begin{proof}
When $s=1$, only $r_+$ is active; when $s=-1$, only $r_-$ is active.
The inactive preactivation is strictly negative by
\eqref{eq:lr-gate-condition}.  Substitution gives the stated identity.
\end{proof}

\paragraph{Rule covariance and bounded perturbations.}
For a purely centered hidden coordinate $H_t=\rho sZ_t$,
Equations~\eqref{eq:lr-innovation-cov} and
\eqref{eq:lr-centered-value} give the exact identity
\begin{equation}\label{eq:lr-centered-seed-cov}
 \mathbb E[g_tH_t]=\rho\chi\bar A(e_1-e_0),\qquad
 \mathbb E[(q_{s_{t-1}}-u)H_t]=0.
\end{equation}
In particular the feature covariance is nonzero even when every
letter site has the same initial attention weight.  If a planted seed
has rule covariance at least $\kappa>0$ and the covariance map is
$L$-Lipschitz on a neighbourhood, then every perturbation of norm at
most $\kappa/(2L)$ preserves covariance at least $\kappa/2$.
This elementary robustness statement applies to perturbations within
a parameterization preserving the feature roles; it does not ensure
that arbitrary off-family trainable directions remain stable under
subsequent unrestricted gradient flow.

\paragraph{Initial positional routing acceleration (ideal centered channel).}
If the hidden representation contains the \emph{exact} independent
coordinate $H_t=\rho sZ_t$, all its other target-correlated coordinates
are independent of positional QK, and $V(0)=0$, then
$F(\phi)=\LRnorm{M(\phi)}_F^2/2$ satisfies, upon varying positional QK,
$F=F_0+\rho^2\chi^2\bar A^2$.  Thus under unit QK mobility,
\begin{equation}\label{eq:lr-centered-accel}
 \dot{\bar A}(0)=0,\qquad
 \ddot{\bar A}(0)=2\rho^2\chi^2\bar A(0)
 \LRnorm{\nabla_{Q,K}\bar A(0)}_F^2\ge0.
\end{equation}
Nondegenerate positional QK gives strict initial acceleration.
This statement is \emph{conditional on exact feature orthogonalization}:
the raw two-ReLU coordinate
$s(d+wZ_t)$ also contains prior information, which creates additional
terms.  It is therefore not a general acceleration theorem for arbitrary
planted Block 2 parameters.

The two hidden units may be trained as an \emph{explicitly tied Euclidean
submanifold}: opposite count weights, equal centered-letter weights and
biases, and opposite output weights.  The induced Euclidean metric assigns
factor $1/2$ to a gradient update of a scalar appearing twice with equal
magnitude.  This tied family is a model restriction; symmetry alone does
not prove its invariance for all unrestricted Transformer parameters.
The mobilities below define flow in the displayed effective factor coordinates.
If a matrix realization uses a fixed direction of nonunit norm, its norm
factor must be absorbed into the assigned coordinate mobility; the scalar
equations are not a claim about unnormalized full-matrix Euclidean flow.

Take $u(0)=c(0)=0$, $\omega_0>0$, $b_0<0$, $v_0>0$, and
$a_0>|\beta_0|+|b_0v_0|$.  Assume all planted magnitudes and gate margins
are inverse-polynomially bounded.  Initialize the positional attention
scores without a preferred correct position and set all trainable
relative-position biases to zero.  This is a \emph{weak computational seed},
not random discovery of the count--letter interaction.  Unlike the
fixed-sign-ReLU-bias initialization of Assumption~2 in the main paper,
the gates here deliberately respond differently to the displayed count.
Thus Theorem~3 does not automatically give a matched outcome plateau for
this precise planted initialization.

\subsection{Exact cross-entropy routing signs}
\label{app:weak-routing-signs}

Let $\sigma(x)=(1+e^{-x})^{-1}$, $y_t'=1-b_t$ and
\begin{equation}\label{eq:lr-logistic-risk}
 \LRstep=\sum_{t=2}^n\omega_t\,\mathbb E
 \left[\log(1+e^{C+gZ_t})-y_t'(C+gZ_t)\right].
\end{equation}
This is exactly the binary cross-entropy of
\eqref{eq:lr-feature-param} because multiplying the binary logit and
label by $s$ cancels the previous-state sign.
Let $A_t=a_{t,j_t}$ and $\bar A=\sum_t\omega_t A_t$.

\begin{lrlemma}[Innovation selects the current source]
\label{lem:lr-signs}
For position-only attention,
\begin{equation}\label{eq:lr-initial-gate-gradient}
 \left.\frac{\partial\LRstep}{\partial g}\right|_{g=0}
 =\chi\bar A>0
\end{equation}
for every $C$.  At any finite logits and $g<0$, define
$v_{tj}=\partial L_t/\partial a_{tj}$, where $L_t$ is the single-time
population risk without $\omega_t$.  Then
\begin{equation}\label{eq:lr-feature-source-signs}
 v_{t,j_t}<0,\qquad v_{tj}\ge0
 \quad(j\in\mathcal J,\;j\ne j_t),\qquad
 v_{tj}=0\quad(j\notin\mathcal J).
\end{equation}
If $|C|+|g|\le B$, then
\begin{equation}\label{eq:lr-source-bounds}
 v_{t,j_t}\le-2\chi|g|\sigma(-B),\qquad
 0\le v_{tj}\le\frac{\chi g^2a_{tj}}4\quad(j\ne j_t\text{ a letter}).
\end{equation}
\end{lrlemma}
\begin{proof}
Centering and independence give
$\mathbb E[(b_t-p)Z_t]=\chi A_t$, yielding
\eqref{eq:lr-initial-gate-gradient}.  For the correct letter,
conditioning on the other bits gives
\[
 v_{t,j_t}=g\chi\mathbb E
 \{1+\sigma(U+gA_t(1-p))-\sigma(U-gA_tp)\},
\]
where $U$ is the contribution of all other positions and $C$.
The braced term is strictly positive.  Under the logit bound it is at
least $2\sigma(-B)$.  For an incorrect letter $j$, the target is
independent of $b_j$, and conditioning on all other bits yields
\[
 v_{tj}=g\chi\mathbb E
 \{\sigma(U+ga_{tj}(1-p))-\sigma(U-ga_{tj}p)\}.
\]
For $g<0$ the two terms in braces are nonpositive; their magnitude is
at most $|g|a_{tj}/4$ by $\sigma'\le1/4$.  Non-letter values are zero.
\end{proof}

Put $G_{tj}=\omega_ta_{tj}(v_{tj}-\bar v_t)$, with
$\bar v_t=\sum_ka_{tk}v_{tk}$; this is the gradient with respect to a
pre-softmax score.  The following global comparison will be important.

\begin{lrlemma}[Correct-versus-incorrect score comparison]
\label{lem:lr-postrouting-gradient}
Whenever $g<0$ and $A_t\ge1/2$,
\begin{equation}\label{eq:lr-score-comparison}
 G_{t,j_t}<0,\qquad G_{tj}-G_{t,j_t}\ge0
 \quad\text{for every other visible site }j.
\end{equation}
Consequently, if the attention biases are shared by relative offset and
$i_t-j_t=n$, then for every $\ell\ne n$,
\begin{equation}\label{eq:lr-bias-gap-mono}
 \frac{d}{dt}
 (\beta_n^{\mathrm{rel}}-\beta_\ell^{\mathrm{rel}})\ge0
\end{equation}
under common positive relative-bias mobility.
\end{lrlemma}
\begin{proof}
The first inequality follows because the correct attention-weight
derivative is strictly smaller than all others, so it lies below their
weighted mean.  For an incorrect site with $a=a_{tj}\le A=A_t$ and
$v=v_{tj}\ge0$, $v_c=v_{t,j_t}<0$,
\[
 \frac{G_{tj}-G_{t,j_t}}{\omega_t}
 =av-Av_c+(A-a)\bar v_t
 \ge av+A(1-A+a)(-v_c)\ge0.
\]
We used $\bar v_t\ge Av_c$ and $A-a\ge0$.
An offset $\ell$ occurs at most once in any attention row.  The correct
offset $n$ appears in every supervised row.  Summing the comparisons
within rows where $\ell$ occurs, and using $G_{t,j_t}<0$ in the
remaining rows, gives
$\partial_{\beta_n^{\mathrm{rel}}}\LRstep
\le\partial_{\beta_\ell^{\mathrm{rel}}}\LRstep$.
\end{proof}

\subsection{Feature bootstrap with every parameter group active}
\label{app:weak-routing-bootstrap}

All training uses one \emph{fixed} system of positive mobilities:
$c,u$ have mobility one; tied MLP/value factors have mobility
$\nu_F>0$; positional QK have mobility $\nu_Q>0$; all relative biases
have common mobility $\nu_R>0$.
No parameter group is explicitly frozen or switched at a chosen time.
The factorized feature family and two-ReLU gate condition are nevertheless
restrictions on the parameterization.

Write $\alpha=u\omega$, $w=bv<0$, $d=a+\beta>0$,
$L_C=\partial\LRstep/\partial C$, $L_g=\partial\LRstep/\partial g$.
While the gate regime holds, the tied Euclidean gradient equations include
\begin{align}
 \dot c&=-L_C,
 &\dot u&=-\omega(dL_C+wL_g),\label{eq:lr-factor-u}\\
 \dot\omega&=-\tfrac{\nu_Fu}{2}(dL_C+wL_g),
 &\dot a=\dot\beta&=-\tfrac{\nu_F\alpha}{2}L_C,\\
 \dot b&=-\tfrac{\nu_F\alpha v}{2}L_g,
 &\dot v&=-\nu_F\alpha bL_g.
 \label{eq:lr-factor-bv}
\end{align}
In particular,
\begin{equation}\label{eq:lr-factor-invariants}
 u^2-\frac{2}{\nu_F}\omega^2=-\frac{2}{\nu_F}\omega_0^2,
 \qquad
 v^2-2b^2=v_0^2-2b_0^2,
 \qquad a-\beta=a_0-\beta_0.
\end{equation}
The first invariant gives $\omega^2\ge\omega_0^2$.  Under the gate
condition, the independent prior readout satisfies
$c(t)\le c_\star=\log((1-p)/p)$: at $c=c_\star$, with $u\omega\ge0$,
its derivative points inward.  At $u=0$ the velocity points into $u>0$,
and $w=bv$ retains its negative sign as long as the bounded gate regime
holds: on either boundary $b=0$ or $v=0$, one has $L_g=\chi\bar A>0$,
so the vector field points into $b<0,v>0$ (the simultaneous zero is excluded
by uniqueness from the nonzero seed). Thus $g=u\omega w<0$ for positive times.

\begin{lrlemma}[Quantitative weak-feature bootstrap, with explicit exits]
\label{lem:lr-feature-bootstrap}
Suppose on an interval prior to reaching the target magnitude that
\begin{equation}\label{eq:lr-amplification-region}
 \bar A\ge A_{\min}>0,\quad
 d\le D_F,\quad
 0<W\le|w|\le W_{\max},\quad
 \omega\ge\omega_{\min}>0,
\end{equation}
and the gate condition \eqref{eq:lr-gate-condition} holds.  Put
\begin{equation}\label{eq:lr-gamma-bootstrap}
 H=D_F+W_{\max},\quad
 \gamma=\chi A_{\min}\left(\frac W H\right)^2,
 \quad
 T_{\mathrm{boot}}=\frac4{3\omega_{\min}^2H^2}.
\end{equation}
Starting with $c=u=0$, the quantity $x=-g$ reaches $x\ge\gamma$
by $T_{\mathrm{boot}}$, unless one of the region conditions fails first.
No assumption $L_C\le0$ is required.
\end{lrlemma}
\begin{proof}
At $c\le c_\star$, let $\alpha=u\omega$.  Since
$0<d+wZ_t\le H$, logistic curvature implies
$L_C\le\alpha H/4$.  At $g=0$, $L_g=\chi\bar A$ and the
same curvature bound gives $L_g\ge\chi A_{\min}-x/4$.
When $x\le\gamma$,
\[
 dL_C+wL_g
 \le\frac{D_FH}{4W}x-W\chi A_{\min}
 +\frac{W_{\max}}4x
 \le-\frac34W\chi A_{\min}.
\]
The last step uses
$D_FH+WW_{\max}\le H^2$.
From \eqref{eq:lr-factor-u} and
$\omega^2\ge\omega_{\min}^2$,
$\dot\alpha\ge(3/4)\omega_{\min}^2W\chi A_{\min}$.
In this region $L_g>0$ and
$\dot w=-\nu_F\alpha(v^2+b^2/2)L_g\le0$.
Hence $\dot x\ge(3/4)\omega_{\min}^2W^2\chi A_{\min}$,
which gives the claimed hitting time.
\end{proof}

The groupwise loss budget is exact in the induced Euclidean metrics:
for the tied feature coordinates $(a,\beta,b,\omega,v)$, the squared
metric weights are $(2,2,2,2,1)$. For projected positional rows $(Q,K)$,
the squared metric weight is $\varepsilon^{-2}$.
\begin{equation}\label{eq:lr-mobility-budget}
 -\dot{\LRstep}=
 \LRnorm{\nabla_{c,u}\LRstep}^2
 +\nu_F\LRnorm{\nabla_F\LRstep}^2
 +\nu_Q\LRnorm{\nabla_{Q,K}\LRstep}^2
 +\nu_R\LRnorm{\nabla_{\beta^{\mathrm{rel}}}\LRstep}^2.
\end{equation}
Here group norms and gradients use those induced metrics. For each group
with mobility $\nu$, Cauchy--Schwarz implies
\begin{equation}\label{eq:lr-movement-budget}
 \LRnorm{\theta_G(T)-\theta_G(0)}
 \le\sqrt{\nu T\,[\LRstep(0)-\LRstep(T)]}
 \le\sqrt{\nu T\log2}.
\end{equation}
Thus sufficiently small but positive $\nu_F,\nu_Q$ preserve prescribed
strict feature-gate and base-score margins over any known finite horizon.
This justifies a controlled bootstrap; it does not establish unit-rate
stability in the unrestricted architecture.

\subsection{Relative-position learning and post-routing convergence}
\label{app:weak-routing-convergence}

Use positional scores
\begin{equation}\label{eq:lr-relative-scores}
 S_{i_tj}=B_{i_tj}+\beta_{i_t-j}^{\mathrm{rel}},
 \qquad
 B_{i_tj}=\frac{q_{i_t}^\top k_j}{\sqrt{d_h}},
 \qquad \beta_\ell^{\mathrm{rel}}(0)=0.
\end{equation}
The correct offset $n$ is not privileged at initialization.

\begin{lrlemma}[Correct-offset learning in the small-feature region]
\label{lem:lr-bias-learning}
Assume all base scores satisfy $|B_{i_tj}|\le D_A$,
all relative biases start at zero, and
$0<x=-g\le A_{\min}/4$. If $A_t\ge A_{\min}>0$,
then the correct score derivative is negative and all other visible-site
score derivatives are nonnegative:
\begin{equation}\label{eq:lr-smallscore-sign}
 G_{t,j_t}<0,\qquad G_{tj}\ge0\ (j\ne j_t).
\end{equation}
Consequently, for the shared relative biases,
\begin{equation}\label{eq:lr-beta-routes}
 \dot\beta_n^{\mathrm{rel}}>0,\qquad
 \dot\beta_\ell^{\mathrm{rel}}\le0\quad(\ell\ne n).
\end{equation}
They imply the uniform floor
\begin{equation}\label{eq:lr-floor}
 A_t\ge A_{\min}:=\frac{e^{-2D_A}}m,
 \qquad m:=\max_t\#\{j:j\le i_t\}\le3n.
\end{equation}
If the initial bootstrap reaches $x=\gamma$ and
$\gamma\le A_{\min}/4$, there is a polynomially bounded fixed bias
mobility $\nu_R$ that drives every $A_t$ to $2/3$ before $x$ leaves
$[\gamma/2,A_{\min}/4]$, provided the feature and base-score
neighbourhoods hold over that interval.
\end{lrlemma}
\begin{proof}
By Lemma~\ref{lem:lr-signs}, the correct value derivative is negative;
incorrect-letter derivatives are nonnegative.  The correct derivative has the logit-independent bound
\[
 -v_{t,j_t}\ge\chi x(1-xA_t/4)\ge\chi x(1-x/4),
\]
because $\sigma$ is $1/4$-Lipschitz. Hence
\[
 \bar v_t\le-\chi x A_t(1-x/4)+\chi x^2/4<0
\]
when $x\le A_{\min}/4$ and $A_t\ge A_{\min}$.  Hence the softmax derivatives have the stated signs.
The bias sign conclusion follows from the fixed offset
\eqref{eq:lr-constant-lag}.
With $\beta_n^{\mathrm{rel}}\ge0$, all other biases $\le0$,
and $|B|\le D_A$, softmax normalization yields
$A_t\ge e^{-2D_A}/m$.

For the time estimate, whenever some row has $A_t\le2/3$, the exact
correct-score derivative satisfies
$-G_{t,j_t}\ge c_0\omega_t\chi xA_{\min}$
with $c_0=1/4$ in this small-$x$ region: indeed
$-G_{t,j_t}\ge\omega_t A_t(1-A_t)(-v_{t,j_t})$.
Therefore
$\dot\beta_n^{\mathrm{rel}}
\ge c_0\nu_R\chi\gamma A_{\min}/[2(n-1)]$
as long as $x\ge\gamma/2$ and routing is incomplete.
All rows achieve $A_t\ge2/3$ when
$\beta_n^{\mathrm{rel}}\ge
2D_A+\log(2(m-1))$.
Thus the additional time is at most
\[
 T_{\mathrm{route}}
 \le \frac{2(n-1)[2D_A+\log(2(m-1))]}
 {c_0\nu_R\chi\gamma A_{\min}}.
\]
On a bounded feature neighbourhood, $|\dot x|\le M_x$ with $M_x$
independent of $\nu_R$ (attention probabilities stay in $[0,1]$).
Choose the fixed $\nu_R$ large enough that this time is less than
$\min\{1,\gamma/(4M_x),(A_{\min}/4-\gamma)/(2M_x)\}$.
For $p=1/4$, the choice $\gamma=\chi A_{\min}(W/H)^2$
automatically gives $\gamma\le3A_{\min}/16<A_{\min}/4$.
No learning-rate switch is required.
\end{proof}

\begin{lrlemma}[Post-routing invariance under bounded QK-score drift]
\label{lem:lr-post-routing}
Suppose $g<0$, $A_t(t_0)\ge2/3$ for all supervised $t$, and the
base-score contrasts change by at most $\log(5/4)$ after $t_0$:
\[
 |[B_{i_t,j_t}-B_{i_tj}](t)-[B_{i_t,j_t}-B_{i_tj}](t_0)|
 \le\log(5/4).
\]
As long as $g<0$ and the centered-feature family holds, all later
correct-position attention weights obey
\begin{equation}\label{eq:lr-post-floor}
 A_t(t)\ge\frac8{13}>\frac35.
\end{equation}
\end{lrlemma}
\begin{proof}
As long as all $A_t\ge1/2$, the bias comparisons of
Lemma~\ref{lem:lr-postrouting-gradient} imply that each relative-bias
gap $\beta_n^{\mathrm{rel}}-\beta_\ell^{\mathrm{rel}}$ is
nondecreasing.  The assumed bound on the other score contrasts says
the odds of each competitor relative to the correct site can increase
by at most a factor $5/4$.  Since the sum of competitor odds is at most
$1/2$ at time $t_0$, it stays at most $5/8$, yielding
$A_t\ge1/(1+5/8)=8/13$.  A first-exit argument closes the condition
$A_t\ge1/2$.
\end{proof}

\begin{lrlemma}[Uniform separability and convex-readout descent]
\label{lem:lr-readout-convergence}
Suppose $A_t\ge3/5$ at every supervised position.  Then
\begin{equation}\label{eq:lr-bit-separation}
 (1-2b_t)\left(\frac12-p-Z_t\right)\ge\frac1{10}
\end{equation}
for every prompt-bit realization.  Suppose in addition that, on a
finite interval after $t_0$, the trainable feature factors remain in a
neighbourhood that admits a \emph{fixed} readout comparator
$\widehat r=(\widehat c,\widehat u)$ satisfying
\begin{equation}\label{eq:lr-comparator}
 \LRstep(\widehat r;\phi(t))\le\eta/2
 \quad\text{for all }t\text{ in the interval}.
\end{equation}
If the two readout coefficients $(c,u)$ continue to train at unit
mobility and every other group has nonnegative mobility, then
\begin{equation}\label{eq:lr-descent}
 \LRstep(t_0+T)\le\frac\eta2+
 \frac{\LRnorm{r(t_0)-\widehat r}^2}{2T}.
\end{equation}
In particular, $T=\LRnorm{r(t_0)-\widehat r}^2/\eta$ suffices for
$\LRstep\le\eta$.
\end{lrlemma}
\begin{proof}
For $b_t=0$, $Z_t\le1-p-A_t$; for $b_t=1$,
$Z_t\ge A_t-p$.  Both cases imply
\eqref{eq:lr-bit-separation}.
For a reference seed $(d_0,w_0,\omega_0)$ with $w_0<0$, take
\[
 R=20\log(2/\eta),\qquad
 \widehat u=\frac{R}{\omega_0|w_0|},\qquad
 \widehat c=R\left(\frac12-p-\frac{d_0}{|w_0|}\right).
\]
At the reference feature, its effective logit is
$R(1/2-p-Z_t)$, which has signed margin at least $R/10$.
Require the feature neighbourhood to satisfy
\begin{equation}\label{eq:lr-comparator-stability}
 |\omega d-\omega_0d_0|+|\omega w-\omega_0w_0|
 \le\frac{\omega_0|w_0|}{20}.
\end{equation}
Since $|Z_t|\le1$, the logit perturbation is at most $R/20$.
The signed margin remains at least $R/20$, so the comparator loss is
at most $e^{-R/20}=\eta/2$. The same comparator is fixed before training,
independently of the routing hitting time.
Cross-entropy is convex in $r=(c,u)$ at fixed feature factors.  Thus
$\frac12\frac{d}{dt}\LRnorm{r-\widehat r}^2
\le\LRstep(\widehat r;\phi(t))-\LRstep(r;\phi(t))$.
The full joint flow makes the loss nonincreasing.  Integrating gives
\eqref{eq:lr-descent}.
\end{proof}

\begin{lrtheorem}[Finite-accuracy learned routing in a symmetry-reduced family]
\label{thm:lr-conditional-bridge}
Fix $n\ge2$, $p=1/4$, $0<\eta<1/2$ and $0<\varepsilon\le1$.
Use the explicitly tied family \eqref{eq:lr-gates}--\eqref{eq:lr-feature-param}
with $c_0=u_0=0$, $b_0<0$, $v_0,\omega_0>0$, and
$a_0>|\beta_0|+|b_0v_0|$. The fixed query and centered value direction
are as in \cref{app:weak-routing-family}. QK consists only of trainable
positional projections of orthogonal position embeddings; every visible
causal site participates in softmax. All relative biases start at zero.

Assume a complexity bound $P\ge2$ bounds $n$, the widths,
$\varepsilon^{-1}$, all initial factor magnitudes, and the Frobenius norms
of the projected QK matrices. Also require
\[
 \min\{|b_0|,v_0,\omega_0,a_0-|\beta_0|-|b_0v_0|\}\ge P^{-1}.
\]
The initial base scores must have moderate spread: for some $D_0\ge0$,
\[
 |B_{i_tj}(0)|\le D_0,\qquad e^{2D_0}\le P.
\]
No correct relative offset is privileged by the relative biases. The base
scores may be bounded perturbations of equal scores; the assertion does
not require an initially favorable correct-site score.

There exist fixed positive mobilities $\nu_F,\nu_Q$ and $\nu_R$, with
$\nu_F^{-1},\nu_Q^{-1},\nu_R$ bounded by polynomials in $P$ and $1/\eta$,
such that the uninterrupted joint population gradient flow reaches
\begin{equation}\label{eq:lr-finite-theorem}
 \LRstep(T)\le\eta
\end{equation}
for $T\le\operatorname{poly}(P,1/\eta)$.
Before the post-routing convergence interval it reaches $A_t\ge2/3$
for every supervised $t$, and preserves $A_t\ge8/13$ throughout that
interval. The mobilities are chosen once at initialization, as below.

The claim concerns a \emph{target-dependent finite-accuracy schedule of
constant group mobilities}, not asymptotic loss convergence under one
fixed mobility choice independent of $\eta$, and not the unrestricted
Transformer parameterization.
\end{lrtheorem}
\begin{proof}[Proof with mobilities fixed before training]
Choose feature and QK balls of radii $r_F,r_Q>0$ around their initial
coordinates. Their inverse radii can be polynomially bounded in $P$.
For the feature ball require the gate condition, $b<0$, $v,\omega>0$,
\eqref{eq:lr-comparator-stability}, and uniform bounds
\[
 0<W\le|w|\le W_{\max},\quad d\le D_F,\quad
 \omega_{\min}\le\omega\le\Omega.
\]
The quantities $W^{-1},\omega_{\min}^{-1},W_{\max},D_F,\Omega$ are
polynomial in $P$. For the QK ball require $|B|\le D_A=D_0+1$ and
\[
 |(B_{i_t,j_t}-B_{i_tj})-(B_{i_t,j_t}(0)-B_{i_tj}(0))|
 \le\tfrac12\log(5/4).
\]
Polynomial maps on bounded seed balls are Lipschitz with polynomial
constants; the assumed inverse-polynomial gate and factor margins give
these radii. Both QK points at any two times in the ball therefore have
contrast differences at most $\log(5/4)$.

Set $A_{\min}=e^{-2D_A}/m$, $H=D_F+W_{\max}$,
$\gamma=\chi A_{\min}(W/H)^2$ and $T_b=T_{\mathrm{boot}}$ from
\eqref{eq:lr-gamma-bootstrap}. These quantities and their relevant
inverses are polynomial in $P$, and $\gamma\le3A_{\min}/16$.
Set $T_{\mathrm{pre}}=T_b+1$. Until that horizon, inside the feature
ball, the unit-mobility readout equations give
\[
 |\dot c|\le1,\qquad |\dot u|\le\Omega H,
 \qquad \|(c,u)\|\le T_{\mathrm{pre}}\sqrt{1+\Omega^2H^2}.
\]
These bounds use only $|\sigma-y'|\le1$ and $|Z_t|\le1$ and do not
depend on bias mobility. With $\nu_F\le1$, the factor equations also
give a polynomial bound $M_x\ge1$ on $|\dot x|$ in this interval,
uniformly over every attention distribution. Choose
\[
 \tau=\min\{1,\gamma/(4M_x),(A_{\min}/4-\gamma)/(2M_x)\},
 \quad K_R=2D_A+\log(2(m-1)),
\]
\[
 \nu_R=\max\left\{1,
 \frac{2(n-1)K_R}{c_0\chi\gamma A_{\min}\tau}\right\},
 \qquad c_0=1/4.
\]
This fixes the relative-bias mobility in advance and makes the routing
bound at most $\tau$. It is polynomially bounded.

Use the fixed reference comparator from the readout lemma and define
\[
 D=T_{\mathrm{pre}}\sqrt{1+\Omega^2H^2}+\|\widehat r\|,
 \qquad T_{\mathrm{post}}=D^2/\eta,\qquad
 T_*=T_{\mathrm{pre}}+T_{\mathrm{post}}.
\]
Now fix the other mobilities by
\[
 \nu_F=\min\{1,r_F^2/(4T_*\log2)\},\qquad
 \nu_Q=\min\{1,r_Q^2/(4T_*\log2)\}.
\]
Their inverses are polynomially bounded. For the tied feature metric,
all squared weights are at least one; for projected QK the weight is
$\varepsilon^{-2}\ge1$. Thus the movement budget bounds ordinary
coordinate movement by $r_F/2$ and $r_Q/2$ up to $T_*$.
A first-exit argument therefore keeps both groups strictly inside their
chosen balls for the whole argument. Finite-time existence follows from
the smooth loss and its gradient-flow energy budget: all mobilities are
finite positive constants, so finite-time parameter escape is excluded.

Inside these balls, $c\le c_\star$ and $u\ge0$ persist by the inward
boundary signs stated above; $\omega>0$ and $w<0$ give $g<0$ for positive
times. In the small-feature region the bias signs imply the attention
floor by a first-exit argument from zero biases. The bootstrap lemma
reaches $x=\gamma$ by $T_b$ unless all routing is already complete;
even in that case the same bootstrap bound applies. Starting at this
hitting time, within at most $\tau$ the bias lemma makes every
$A_t\ge2/3$ (if this already holds, the extra time is zero). Let $t_0$
be the end of this interval; then $t_0\le T_{\mathrm{pre}}$.
The derivative bound keeps $x$ inside $[\gamma/2,A_{\min}/4]$ during
this routing interval. The bounded base-score contrasts and the
post-routing lemma preserve $A_t\ge8/13>3/5$ afterwards.
The feature ball preserves the fixed comparator for the entire remaining
interval. Its initial distance at $t_0$ is at most $D$, so the readout
lemma gives $\LRstep(t_0+T_{\mathrm{post}})\le\eta$.
All mobilities were selected from initial uniform ball bounds; none
requires a later switch or a posterior choice of $t_0$.
\end{proof}

\subsection{Limitations and relationship to the main comparison}
\label{app:weak-routing-limitations}

\paragraph{What is new.}
The statistical signal for the correct relative offset is derived from
independent innovations rather than inserted in the attention mask or bias
initialization.  All visible causal sites compete in the softmax.  The
tied two-ReLU coefficients, value amplitude, positional QK, linear readout,
and relative-position biases train \emph{simultaneously}, although at
different fixed mobilities.  The constant-offset identity
\eqref{eq:lr-constant-lag} is a genuine architectural inductive bias.

\paragraph{What is not proved.}
The reduced family preselects a state--letter feature direction and
centers prompt-letter values.  It excludes unrestricted content-QKV
updates and arbitrary perturbations that leave this tied family.
The result also permits a large relative-bias mobility and very small
positive feature/QK mobilities, and both may depend on the target error
$\eta$.  This is therefore weaker than a theorem for ordinary
unit-mobility, randomly initialized Block 2.  In addition, the
count-dependent ReLU activity differs from the fixed-sign near-zero
bias regime in the main paper's Assumption~2.  The original outcome
plateau (Theorem~3) cannot be transferred to this planted seed without
an additional argument.  Nor does the finite-accuracy guarantee imply
$\LRstep(t)\to0$ for a single mobility choice independent of $\eta$.

\paragraph{Scope of the audited parameterization.}
For interpreting this appendix extension, the tied-subfamily
symmetry, parameter scaling relative to Assumption~2, and value-role restriction must be distinguished
from unrestricted updates in the original computational graph. The quantitative construction above
uses the stated reduced-coordinate metrics, moderate base-score spread,
and bounded seeds; it does not transfer to arbitrary parameter metrics.  In particular, the
centering of the fixed query-letter token must be implemented by the
value projection rather than ignored, the ordinary token-content QK
directions must be either controlled or explicitly excluded, and any
extension to Theorem~3 must separately verify its activation-pattern
hypotheses.  An earlier inference that \emph{nonzero rule covariance by
itself} forces correct routing is false; count-prior interference can
make the initial positional routing acceleration negative.

\paragraph{Suggested interpretation.}
This result is a complementary \emph{conditional routing bridge} next
to the original Theorem~4, not a replacement of its rigorous
positive-margin/global-loss convergence statement.  The next genuinely
stronger theorem would remove the tied feature restriction while
keeping a planted feature, and then establish feature stability,
relative-position learning, and polynomial-time low loss for all
trainable Block 2 directions, with a matched outcome comparison.

\paragraph{Numerical audit.} Exact Bernoulli enumeration and automatic differentiation in
\nolinkurl{handcoded/theory_checks/learned_routing_extension_checks.py}
check the source gradients, small-feature and post-routing score signs, tied-factor
invariants, positional QK metric, uniform separability, and the fixed comparator.
These checks support the displayed local identities and inequalities; they are
not a formal verification of the global theorem or a test of unrestricted training.
